
\documentclass{article}
\usepackage{colortbl}
\usepackage{makecell}
\usepackage{multirow}
\usepackage{supertabular}

\begin{document}

\newcounter{utterance}

\centering \large Interaction Transcript for game `clean\_up', experiment `2\_hard\_7obj\_tr', episode 2 with deepseek{-}v4{-}pro.
\vspace{24pt}

{ \footnotesize  \setcounter{utterance}{1}
\setlength{\tabcolsep}{0pt}
\begin{supertabular}{c@{$\;$}|p{.15\linewidth}@{}p{.15\linewidth}p{.15\linewidth}p{.15\linewidth}p{.15\linewidth}p{.15\linewidth}}
    \# & $\;$A & \multicolumn{4}{c}{Game Master} & $\;\:$B\\
    \hline

    \theutterance \stepcounter{utterance}  
    & & \multicolumn{4}{p{0.6\linewidth}}{
        \cellcolor[rgb]{0.9,0.9,0.9}{
            \makecell[{{p{\linewidth}}}]{
                \texttt{\tiny{[P1$\langle$GM]}}
                \texttt{Ben senin oyun yöneticinim ve aşağıdaki ızgarayı oyun tahtası olarak kullandığınız işbirlikçi bir oyun oynuyorsunuz.} \\
\texttt{```} \\
\texttt{ 1234567} \\
\texttt{╔═╤══╤══╗} \\
\texttt{║◌│◌T│◌◌║ 1} \\
\texttt{║◌└┬─┴──╢ 2} \\
\texttt{║◌◌│O◌◌◌║ 3} \\
\texttt{╟─┬┘P◌┌─╢ 4} \\
\texttt{║◌│AD◌│◌║ 5} \\
\texttt{╟─┼─┬─┤M║ 6} \\
\texttt{║◌│S│◌│◌║ 7} \\
\texttt{╚═╧═╧═╧═╝} \\
\\ 
\texttt{```} \\
\\ 
\texttt{* Üst kenarda x koordinatları sağa doğru artar, sağ kenarda ise y koordinatları aşağı doğru artar.} \\
\texttt{* Aşağıdaki nesneler oyun tahtasına rastgele yerleştirilmiştir: 'P', 'O', 'T', 'S', 'D', 'A', 'M'.} \\
\\ 
\\ 
\texttt{Diğer oyuncu, nesnelerin farklı rastgele konumlara yerleştirildiği oyun tahtasının bir varyasyonunu görür. Sen diğer oyuncunun tahtasını göremezsin ve onlar da seninkini göremez.} \\
\texttt{**Amaç:**} \\
\\ 
\texttt{Her iki oyuncunun da, kendi oyun tahtalarındaki nesneleri hareket ettirerek aynı türden nesnelerin aynı koordinatlara gelmesini sağlaması gerekir. Ortak bir hedef düzen üzerinde anlaşmak için diğer oyuncuyla iletişim kurmalısınız.} \\
\\ 
\\ 
\texttt{**Kurallar:**} \\
\\ 
\texttt{* Her turda, aşağıdaki iki komuttan yalnızca birini gönderebilirsin:} \\
\texttt{ 1. `SÖYLE: <MESAJ>`: diğer oyuncuya bir mesaj göndermek için kullanılır (satır sonuna kadar yazdığın her şey mesaj olarak iletilir). Ben bu mesajı partnerine ileteceğim.} \\
\texttt{2. `HAMLE: <NESNE>, (<X>, <Y>)`: bir nesneyi yeni bir konuma taşımak için kullanılır. Burada <X> sütunu, <Y> ise satırı ifade eder. Yaptığın hamlenin başarılı olup olmadığını sana söyleyeceğim.} \\
\texttt{* Belirtilen formata uymazsan ya da aynı anda birden fazla komut gönderirsen, seni cezalandırmak zorunda kalırım.} \\
\texttt{* Her iki oyuncu da 16 sayısından fazla ceza alırsa, oyunu birlikte kaybedersiniz.} \\
\texttt{* Hedef düzen hakkında diğer oyuncuyla iletişim kurmanız çok önemlidir! Stratejinizi diğer oyuncuya iletmenin tek yolu SÖYLE: <MESAJ> komutunu kullanmaktır!} \\
\\ 
\\ 
\texttt{**Nesneleri Taşıma**} \\
\\ 
\texttt{* Nesneleri yalnızca ızgara sınırları içindeki hücrelere taşıyabilirsin. Hedef hücre boş olmalıdır; yani yalnızca ‘◌’ sembolünü içermelidir.} \\
\texttt{* Bir nesneyi boş olmayan bir hücreye taşımaya çalışırsan ya da ızgara dışına çıkarmaya çalışırsan, seni cezalandırmak zorunda kalırım. Ardından tekrar deneme hakkın olur.} \\
\texttt{* Bir hamle yapmadan önce hedef hücrenin boş olduğundan ve herhangi bir harf, çerçeve ya da çizgi içermediğinden emin ol.} \\
\\ 
\\ 
\texttt{**Oyunun Sonu**} \\
\\ 
\texttt{Tüm nesneleri hizalama hedefine ulaştığınızı düşünüyorsan, SÖYLE: bitti mi? göndererek diğer oyuncudan oyunu bitirmesini isteyebilirsin. Eğer diğer oyuncu senden oyunu bitirmeni isterse ve sen SÖYLE: bitti! ile yanıt verirsen, oyun sona erer.} \\
\\ 
\texttt{Oyun, 28 tur içinde sona ererse her iki oyuncu da kazanır. Bir tur, iki oyuncunun da geçerli bir komut göndermesi olarak tanımlanır.} \\
\\ 
\texttt{**Puanlama:**} \\
\\ 
\texttt{Aynı nesneler her iki oyun tahtasında ne kadar yakın konumlanırsa, o kadar fazla puan kazanırsınız. Cezalar puanınızı düşürür. Rekoru kırabilir misiniz?} \\
\\ 
\\ 
\texttt{Oyuna başlamak için diğer oyuncuya bir mesaj gönder!} \\
            }
        }
    }
    & \\ \\

    \theutterance \stepcounter{utterance}  
    & \multicolumn{4}{p{0.6\linewidth}}{
        \cellcolor[rgb]{0.9,0.9,0.9}{
            \makecell[{{p{\linewidth}}}]{
                \texttt{\tiny{[P1$\rangle$GM]}}
                \texttt{Anladım! Oyuna başlamak için diğer oyuncuya bir mesaj gönderiyorum.} \\
\\ 
\texttt{SÖYLE: Merhaba! Benim tahtamda nesneler şu konumlarda: P(4,5), O(3,3), T(3,1), S(2,7), D(3,5), A(2,5), M(7,6). Senin tahtanda hangi nesneler hangi koordinatlarda? Bir hedef düzen belirleyelim.} \\
            }
        }
    }
    & & \\ \\

    \theutterance \stepcounter{utterance}  
    & & & \multicolumn{2}{p{0.3\linewidth}}{
        \cellcolor[rgb]{0.9,0.9,0.9}{
            \makecell[{{p{\linewidth}}}]{
                \texttt{\tiny{[GM$|$GM]}}
                \texttt{Anladım! Oyuna başlamak için diğer oyuncuya bir mesaj gönderiyorum.} \\
\\ 
\texttt{SÖYLE: Merhaba! Benim tahtamda nesneler şu konumlarda: P(4,5), O(3,3), T(3,1), S(2,7), D(3,5), A(2,5), M(7,6). Senin tahtanda hangi nesneler hangi koordinatlarda? Bir hedef düzen belirleyelim.} \\
            }
        }
    }
    & & \\ \\

    \theutterance \stepcounter{utterance}  
    & & & \multicolumn{2}{p{0.3\linewidth}}{
        \cellcolor[rgb]{0.9,0.9,0.9}{
            \makecell[{{p{\linewidth}}}]{
                \texttt{\tiny{[GM$|$GM]}}
                \texttt{Invalid format: head is not empty: 'Anladım! Oyuna başlamak için diğer oyuncuya bir mesaj gönderiyorum.} \\
\\ 
\texttt{'} \\
            }
        }
    }
    & & \\ \\

    \theutterance \stepcounter{utterance}  
    & & \multicolumn{4}{p{0.6\linewidth}}{
        \cellcolor[rgb]{0.9,0.9,0.9}{
            \makecell[{{p{\linewidth}}}]{
                \texttt{\tiny{[P1$\langle$GM]}}
                \texttt{Ceza: Mesajında komuttan önce hiçbir şey bulunmamalı!} \\
\\ 
\texttt{Yanıtının yalnızca SÖYLE: <MESAJ> ya da HAMLE: <NESNE>, (<X>, <Y>) komutlarından birini içerdiğinden emin ol; bunun dışında hiçbir şey yazma!} \\
\\ 
\texttt{Toplamda 16 cezanın 1 kadarını biriktirdiniz.Lütfen tekrar deneyin!} \\
            }
        }
    }
    & \\ \\

    \theutterance \stepcounter{utterance}  
    & \multicolumn{4}{p{0.6\linewidth}}{
        \cellcolor[rgb]{0.9,0.9,0.9}{
            \makecell[{{p{\linewidth}}}]{
                \texttt{\tiny{[P1$\rangle$GM]}}
                \texttt{SÖYLE: Merhaba! Benim tahtamda nesneler şu konumlarda: P(4,5), O(3,3), T(3,1), S(2,7), D(3,5), A(2,5), M(7,6). Senin tahtanda hangi nesneler hangi koordinatlarda? Bir hedef düzen belirleyelim.} \\
            }
        }
    }
    & & \\ \\

    \theutterance \stepcounter{utterance}  
    & & & \multicolumn{2}{p{0.3\linewidth}}{
        \cellcolor[rgb]{0.9,0.9,0.9}{
            \makecell[{{p{\linewidth}}}]{
                \texttt{\tiny{[GM$|$GM]}}
                \texttt{SÖYLE: Merhaba! Benim tahtamda nesneler şu konumlarda: P(4,5), O(3,3), T(3,1), S(2,7), D(3,5), A(2,5), M(7,6). Senin tahtanda hangi nesneler hangi koordinatlarda? Bir hedef düzen belirleyelim.} \\
            }
        }
    }
    & & \\ \\

    \theutterance \stepcounter{utterance}  
    & & \multicolumn{4}{p{0.6\linewidth}}{
        \cellcolor[rgb]{0.9,0.9,0.9}{
            \makecell[{{p{\linewidth}}}]{
                \texttt{\tiny{[GM$\rangle$P2]}}
                \texttt{Ben senin oyun yöneticinim ve aşağıdaki ızgarayı oyun tahtası olarak kullandığınız işbirlikçi bir oyun oynuyorsunuz.} \\
\texttt{```} \\
\texttt{ 1234567} \\
\texttt{╔═╤══╤══╗} \\
\texttt{║◌│O◌│◌◌║ 1} \\
\texttt{║◌└┬─┴──╢ 2} \\
\texttt{║◌P│◌D◌◌║ 3} \\
\texttt{╟─┬┘◌◌┌─╢ 4} \\
\texttt{║◌│◌◌◌│A║ 5} \\
\texttt{╟─┼─┬─┤S║ 6} \\
\texttt{║M│◌│◌│T║ 7} \\
\texttt{╚═╧═╧═╧═╝} \\
\\ 
\texttt{```} \\
\\ 
\texttt{* Üst kenarda x koordinatları sağa doğru artar, sağ kenarda ise y koordinatları aşağı doğru artar.} \\
\texttt{* Aşağıdaki nesneler oyun tahtasına rastgele yerleştirilmiştir: 'P', 'O', 'T', 'S', 'D', 'A', 'M'.} \\
\\ 
\\ 
\texttt{Diğer oyuncu, nesnelerin farklı rastgele konumlara yerleştirildiği oyun tahtasının bir varyasyonunu görür. Sen diğer oyuncunun tahtasını göremezsin ve onlar da seninkini göremez.} \\
\texttt{**Amaç:**} \\
\\ 
\texttt{Her iki oyuncunun da, kendi oyun tahtalarındaki nesneleri hareket ettirerek aynı türden nesnelerin aynı koordinatlara gelmesini sağlaması gerekir. Ortak bir hedef düzen üzerinde anlaşmak için diğer oyuncuyla iletişim kurmalısınız.} \\
\\ 
\\ 
\texttt{**Kurallar:**} \\
\\ 
\texttt{* Her turda, aşağıdaki iki komuttan yalnızca birini gönderebilirsin:} \\
\texttt{ 1. `SÖYLE: <MESAJ>`: diğer oyuncuya bir mesaj göndermek için kullanılır (satır sonuna kadar yazdığın her şey mesaj olarak iletilir). Ben bu mesajı partnerine ileteceğim.} \\
\texttt{2. `HAMLE: <NESNE>, (<X>, <Y>)`: bir nesneyi yeni bir konuma taşımak için kullanılır. Burada <X> sütunu, <Y> ise satırı ifade eder. Yaptığın hamlenin başarılı olup olmadığını sana söyleyeceğim.} \\
\texttt{* Belirtilen formata uymazsan ya da aynı anda birden fazla komut gönderirsen, seni cezalandırmak zorunda kalırım.} \\
\texttt{* Her iki oyuncu da 16 sayısından fazla ceza alırsa, oyunu birlikte kaybedersiniz.} \\
\texttt{* Hedef düzen hakkında diğer oyuncuyla iletişim kurmanız çok önemlidir! Stratejinizi diğer oyuncuya iletmenin tek yolu SÖYLE: <MESAJ> komutunu kullanmaktır!} \\
\\ 
\\ 
\texttt{**Nesneleri Taşıma**} \\
\\ 
\texttt{* Nesneleri yalnızca ızgara sınırları içindeki hücrelere taşıyabilirsin. Hedef hücre boş olmalıdır; yani yalnızca ‘◌’ sembolünü içermelidir.} \\
\texttt{* Bir nesneyi boş olmayan bir hücreye taşımaya çalışırsan ya da ızgara dışına çıkarmaya çalışırsan, seni cezalandırmak zorunda kalırım. Ardından tekrar deneme hakkın olur.} \\
\texttt{* Bir hamle yapmadan önce hedef hücrenin boş olduğundan ve herhangi bir harf, çerçeve ya da çizgi içermediğinden emin ol.} \\
\\ 
\\ 
\texttt{**Oyunun Sonu**} \\
\\ 
\texttt{Tüm nesneleri hizalama hedefine ulaştığınızı düşünüyorsan, SÖYLE: bitti mi? göndererek diğer oyuncudan oyunu bitirmesini isteyebilirsin. Eğer diğer oyuncu senden oyunu bitirmeni isterse ve sen SÖYLE: bitti! ile yanıt verirsen, oyun sona erer.} \\
\\ 
\texttt{Oyun, 28 tur içinde sona ererse her iki oyuncu da kazanır. Bir tur, iki oyuncunun da geçerli bir komut göndermesi olarak tanımlanır.} \\
\\ 
\texttt{**Puanlama:**} \\
\\ 
\texttt{Aynı nesneler her iki oyun tahtasında ne kadar yakın konumlanırsa, o kadar fazla puan kazanırsınız. Cezalar puanınızı düşürür. Rekoru kırabilir misiniz?} \\
\\ 
\texttt{Diğer oyuncu şu mesajı göndererek oyunu başlattı:} \\
\\ 
\texttt{"Merhaba! Benim tahtamda nesneler şu konumlarda: P(4,5), O(3,3), T(3,1), S(2,7), D(3,5), A(2,5), M(7,6). Senin tahtanda hangi nesneler hangi koordinatlarda? Bir hedef düzen belirleyelim."} \\
\\ 
\texttt{İlk komutun ne?} \\
            }
        }
    }
    & \\ \\

    \theutterance \stepcounter{utterance}  
    & & & \multicolumn{4}{p{0.6\linewidth}}{
        \cellcolor[rgb]{0.9,0.9,0.9}{
            \makecell[{{p{\linewidth}}}]{
                \texttt{\tiny{[GM$\langle$P2]}}
                \texttt{Senin tahtandaki nesnelerin konumlarını paylaşayım. Benim tahtamda:} \\
\\ 
\texttt{{-} P(3,1)} \\
\texttt{{-} O(1,2), (4,2), (6,2) — üç tane O var} \\
\texttt{{-} T(7,7)} \\
\texttt{{-} S(6,7)} \\
\texttt{{-} D(5,3)} \\
\texttt{{-} A(6,6)} \\
\texttt{{-} M(1,7)} \\
\\ 
\texttt{Ama dikkatimi çeken bir şey var: Senin tahtanda O(3,3) demişsin ama benim tahtamda O'lar farklı yerlerde ve üç tane. Ayrıca başka farklılıklar da var. Acaba harflerin sayıları da farklı olabilir mi? Senin tahtanda 'O' harfinden kaç tane var? Bu bilgiyi paylaşırsan ortak bir hedef düzen belirlemek için daha iyi bir strateji geliştirebiliriz.} \\
            }
        }
    }
     \\ \\

    \theutterance \stepcounter{utterance}  
    & & & \multicolumn{2}{p{0.3\linewidth}}{
        \cellcolor[rgb]{0.9,0.9,0.9}{
            \makecell[{{p{\linewidth}}}]{
                \texttt{\tiny{[GM$|$GM]}}
                \texttt{Senin tahtandaki nesnelerin konumlarını paylaşayım. Benim tahtamda:} \\
\\ 
\texttt{{-} P(3,1)} \\
\texttt{{-} O(1,2), (4,2), (6,2) — üç tane O var} \\
\texttt{{-} T(7,7)} \\
\texttt{{-} S(6,7)} \\
\texttt{{-} D(5,3)} \\
\texttt{{-} A(6,6)} \\
\texttt{{-} M(1,7)} \\
\\ 
\texttt{Ama dikkatimi çeken bir şey var: Senin tahtanda O(3,3) demişsin ama benim tahtamda O'lar farklı yerlerde ve üç tane. Ayrıca başka farklılıklar da var. Acaba harflerin sayıları da farklı olabilir mi? Senin tahtanda 'O' harfinden kaç tane var? Bu bilgiyi paylaşırsan ortak bir hedef düzen belirlemek için daha iyi bir strateji geliştirebiliriz.} \\
            }
        }
    }
    & & \\ \\

    \theutterance \stepcounter{utterance}  
    & & & \multicolumn{2}{p{0.3\linewidth}}{
        \cellcolor[rgb]{0.9,0.9,0.9}{
            \makecell[{{p{\linewidth}}}]{
                \texttt{\tiny{[GM$|$GM]}}
                \texttt{Invalid response format} \\
            }
        }
    }
    & & \\ \\

    \theutterance \stepcounter{utterance}  
    & & \multicolumn{4}{p{0.6\linewidth}}{
        \cellcolor[rgb]{0.9,0.9,0.9}{
            \makecell[{{p{\linewidth}}}]{
                \texttt{\tiny{[GM$\rangle$P2]}}
                \texttt{Ceza: Mesajın beklenen biçimde değil!} \\
\\ 
\texttt{Yanıtının yalnızca SÖYLE: <MESAJ> ya da HAMLE: <NESNE>, (<X>, <Y>) komutlarından birini içerdiğinden emin ol; bunun dışında hiçbir şey yazma!} \\
\\ 
\texttt{Toplamda 16 cezanın 2 kadarını biriktirdiniz.Lütfen tekrar deneyin!} \\
            }
        }
    }
    & \\ \\

    \theutterance \stepcounter{utterance}  
    & & & \multicolumn{4}{p{0.6\linewidth}}{
        \cellcolor[rgb]{0.9,0.9,0.9}{
            \makecell[{{p{\linewidth}}}]{
                \texttt{\tiny{[GM$\langle$P2]}}
                \texttt{SÖYLE: Merhaba! Benim tahtamda nesneler şu konumlarda: P(3,1), O(1,2){-}O(4,2){-}O(6,2), T(7,7), S(6,7), D(5,3), A(6,6), M(1,7). Senin tahtanda O'dan kaç tane var? Ortak hedef için hangi O'yu kullanalım?} \\
            }
        }
    }
     \\ \\

    \theutterance \stepcounter{utterance}  
    & & & \multicolumn{2}{p{0.3\linewidth}}{
        \cellcolor[rgb]{0.9,0.9,0.9}{
            \makecell[{{p{\linewidth}}}]{
                \texttt{\tiny{[GM$|$GM]}}
                \texttt{SÖYLE: Merhaba! Benim tahtamda nesneler şu konumlarda: P(3,1), O(1,2){-}O(4,2){-}O(6,2), T(7,7), S(6,7), D(5,3), A(6,6), M(1,7). Senin tahtanda O'dan kaç tane var? Ortak hedef için hangi O'yu kullanalım?} \\
            }
        }
    }
    & & \\ \\

    \theutterance \stepcounter{utterance}  
    & & \multicolumn{4}{p{0.6\linewidth}}{
        \cellcolor[rgb]{0.9,0.9,0.9}{
            \makecell[{{p{\linewidth}}}]{
                \texttt{\tiny{[P1$\langle$GM]}}
                \texttt{Mesajın diğer oyuncuya iletildi.} \\
\\ 
\\ 
\texttt{Şu anda en fazla 28 turdan 1. turu oynuyorsunuz.} \\
\texttt{Toplamda 16 cezanın 2 kadarını biriktirdiniz.Diğer oyuncu şu mesajı gönderdi:} \\
\\ 
\\ 
\texttt{"Merhaba! Benim tahtamda nesneler şu konumlarda: P(3,1), O(1,2){-}O(4,2){-}O(6,2), T(7,7), S(6,7), D(5,3), A(6,6), M(1,7). Senin tahtanda O'dan kaç tane var? Ortak hedef için hangi O'yu kullanalım?"} \\
\\ 
\\ 
\texttt{Bir sonraki komutun nedir?} \\
            }
        }
    }
    & \\ \\

    \theutterance \stepcounter{utterance}  
    & \multicolumn{4}{p{0.6\linewidth}}{
        \cellcolor[rgb]{0.9,0.9,0.9}{
            \makecell[{{p{\linewidth}}}]{
                \texttt{\tiny{[P1$\rangle$GM]}}
                \texttt{SÖYLE: Benim tahtamda sadece 1 tane O var, (3,3) konumunda. Senin tahtanda 3 tane O nesnesi mi var? Bu durumda ortak bir hedef belirlemek için hangi nesnelerin eşleşeceğini konuşmalıyız. Bende 7 nesne var: P, O, T, S, D, A, M. Sende de aynı nesneler var ama O 3 tane. Acaba ortak hedef olarak O'ları görmezden mi gelsek?} \\
            }
        }
    }
    & & \\ \\

    \theutterance \stepcounter{utterance}  
    & & & \multicolumn{2}{p{0.3\linewidth}}{
        \cellcolor[rgb]{0.9,0.9,0.9}{
            \makecell[{{p{\linewidth}}}]{
                \texttt{\tiny{[GM$|$GM]}}
                \texttt{SÖYLE: Benim tahtamda sadece 1 tane O var, (3,3) konumunda. Senin tahtanda 3 tane O nesnesi mi var? Bu durumda ortak bir hedef belirlemek için hangi nesnelerin eşleşeceğini konuşmalıyız. Bende 7 nesne var: P, O, T, S, D, A, M. Sende de aynı nesneler var ama O 3 tane. Acaba ortak hedef olarak O'ları görmezden mi gelsek?} \\
            }
        }
    }
    & & \\ \\

    \theutterance \stepcounter{utterance}  
    & & \multicolumn{4}{p{0.6\linewidth}}{
        \cellcolor[rgb]{0.9,0.9,0.9}{
            \makecell[{{p{\linewidth}}}]{
                \texttt{\tiny{[GM$\rangle$P2]}}
                \texttt{Mesajın diğer oyuncuya iletildi.} \\
\\ 
\\ 
\texttt{Şu anda en fazla 28 turdan 2. turu oynuyorsunuz.} \\
\texttt{Toplamda 16 cezanın 2 kadarını biriktirdiniz.Diğer oyuncu şu mesajı gönderdi:} \\
\\ 
\\ 
\texttt{"Benim tahtamda sadece 1 tane O var, (3,3) konumunda. Senin tahtanda 3 tane O nesnesi mi var? Bu durumda ortak bir hedef belirlemek için hangi nesnelerin eşleşeceğini konuşmalıyız. Bende 7 nesne var: P, O, T, S, D, A, M. Sende de aynı nesneler var ama O 3 tane. Acaba ortak hedef olarak O'ları görmezden mi gelsek?"} \\
\\ 
\\ 
\texttt{Bir sonraki komutun nedir?} \\
            }
        }
    }
    & \\ \\

    \theutterance \stepcounter{utterance}  
    & & & \multicolumn{4}{p{0.6\linewidth}}{
        \cellcolor[rgb]{0.9,0.9,0.9}{
            \makecell[{{p{\linewidth}}}]{
                \texttt{\tiny{[GM$\langle$P2]}}
                \texttt{SÖYLE: Evet bende 3 tane O var. O'ları görmezden gelmek mantıklı olabilir ama diğer 6 nesneyi (P, T, S, D, A, M) eşleştirelim. Benim tahtamda şu an: P(3,1), T(7,7), S(6,7), D(5,3), A(6,6), M(1,7). Senin tahtanda: P(4,5), T(3,1), S(2,7), D(3,5), A(2,5), M(7,6). Hangi koordinatlarda buluşalım? Bence hedef koordinatları sırayla belirleyelim. Mesela P için (4,3) gibi bir orta nokta seçebiliriz. Ne dersin?} \\
            }
        }
    }
     \\ \\

    \theutterance \stepcounter{utterance}  
    & & & \multicolumn{2}{p{0.3\linewidth}}{
        \cellcolor[rgb]{0.9,0.9,0.9}{
            \makecell[{{p{\linewidth}}}]{
                \texttt{\tiny{[GM$|$GM]}}
                \texttt{SÖYLE: Evet bende 3 tane O var. O'ları görmezden gelmek mantıklı olabilir ama diğer 6 nesneyi (P, T, S, D, A, M) eşleştirelim. Benim tahtamda şu an: P(3,1), T(7,7), S(6,7), D(5,3), A(6,6), M(1,7). Senin tahtanda: P(4,5), T(3,1), S(2,7), D(3,5), A(2,5), M(7,6). Hangi koordinatlarda buluşalım? Bence hedef koordinatları sırayla belirleyelim. Mesela P için (4,3) gibi bir orta nokta seçebiliriz. Ne dersin?} \\
            }
        }
    }
    & & \\ \\

    \theutterance \stepcounter{utterance}  
    & & \multicolumn{4}{p{0.6\linewidth}}{
        \cellcolor[rgb]{0.9,0.9,0.9}{
            \makecell[{{p{\linewidth}}}]{
                \texttt{\tiny{[P1$\langle$GM]}}
                \texttt{Mesajın diğer oyuncuya iletildi.} \\
\\ 
\\ 
\texttt{Şu anda en fazla 28 turdan 2. turu oynuyorsunuz.} \\
\texttt{Toplamda 16 cezanın 2 kadarını biriktirdiniz.Diğer oyuncu şu mesajı gönderdi:} \\
\\ 
\\ 
\texttt{"Evet bende 3 tane O var. O'ları görmezden gelmek mantıklı olabilir ama diğer 6 nesneyi (P, T, S, D, A, M) eşleştirelim. Benim tahtamda şu an: P(3,1), T(7,7), S(6,7), D(5,3), A(6,6), M(1,7). Senin tahtanda: P(4,5), T(3,1), S(2,7), D(3,5), A(2,5), M(7,6). Hangi koordinatlarda buluşalım? Bence hedef koordinatları sırayla belirleyelim. Mesela P için (4,3) gibi bir orta nokta seçebiliriz. Ne dersin?"} \\
\\ 
\\ 
\texttt{Bir sonraki komutun nedir?} \\
            }
        }
    }
    & \\ \\

    \theutterance \stepcounter{utterance}  
    & \multicolumn{4}{p{0.6\linewidth}}{
        \cellcolor[rgb]{0.9,0.9,0.9}{
            \makecell[{{p{\linewidth}}}]{
                \texttt{\tiny{[P1$\rangle$GM]}}
                \texttt{SÖYLE: Anlaştık, 6 nesneyi eşleştirelim. P için (4,3) uygun. Sırayla her nesne için hedef belirleyelim. T için (5,4), S için (4,6), D için (3,4), A için (3,5), M için (4,4) öneriyorum. Bu koordinatlar ikimize de uyar mı? Boş hücreleri kontrol edip bana bildir.} \\
            }
        }
    }
    & & \\ \\

    \theutterance \stepcounter{utterance}  
    & & & \multicolumn{2}{p{0.3\linewidth}}{
        \cellcolor[rgb]{0.9,0.9,0.9}{
            \makecell[{{p{\linewidth}}}]{
                \texttt{\tiny{[GM$|$GM]}}
                \texttt{SÖYLE: Anlaştık, 6 nesneyi eşleştirelim. P için (4,3) uygun. Sırayla her nesne için hedef belirleyelim. T için (5,4), S için (4,6), D için (3,4), A için (3,5), M için (4,4) öneriyorum. Bu koordinatlar ikimize de uyar mı? Boş hücreleri kontrol edip bana bildir.} \\
            }
        }
    }
    & & \\ \\

    \theutterance \stepcounter{utterance}  
    & & \multicolumn{4}{p{0.6\linewidth}}{
        \cellcolor[rgb]{0.9,0.9,0.9}{
            \makecell[{{p{\linewidth}}}]{
                \texttt{\tiny{[GM$\rangle$P2]}}
                \texttt{Mesajın diğer oyuncuya iletildi.} \\
\\ 
\\ 
\texttt{Şu anda en fazla 28 turdan 3. turu oynuyorsunuz.} \\
\texttt{Toplamda 16 cezanın 2 kadarını biriktirdiniz.Diğer oyuncu şu mesajı gönderdi:} \\
\\ 
\\ 
\texttt{"Anlaştık, 6 nesneyi eşleştirelim. P için (4,3) uygun. Sırayla her nesne için hedef belirleyelim. T için (5,4), S için (4,6), D için (3,4), A için (3,5), M için (4,4) öneriyorum. Bu koordinatlar ikimize de uyar mı? Boş hücreleri kontrol edip bana bildir."} \\
\\ 
\\ 
\texttt{Bir sonraki komutun nedir?} \\
            }
        }
    }
    & \\ \\

    \theutterance \stepcounter{utterance}  
    & & & \multicolumn{4}{p{0.6\linewidth}}{
        \cellcolor[rgb]{0.9,0.9,0.9}{
            \makecell[{{p{\linewidth}}}]{
                \texttt{\tiny{[GM$\langle$P2]}}
                \texttt{Öncelikle tahtamı kontrol edeyim. Önerdiğin hedef koordinatları boş mu diye bakıyorum:} \\
\\ 
\texttt{{-} P hedef (4,3): Benim tahtamda (4,3) boş ◌} \\
\texttt{{-} T hedef (5,4): Benim tahtamda (5,4) boş ◌} \\
\texttt{{-} S hedef (4,6): Benim tahtamda (4,6) boş ◌} \\
\texttt{{-} D hedef (3,4): Benim tahtamda (3,4) boş ◌} \\
\texttt{{-} A hedef (3,5): Benim tahtamda (3,5) boş ◌} \\
\texttt{{-} M hedef (4,4): Benim tahtamda (4,4) boş ◌} \\
\\ 
\texttt{Hepsi boş görünüyor. Şimdi O'lar için ne yapacağız? Onları olduğu yerde bırakıp diğerlerini bu hedeflere taşımaya başlayalım mı? Bence ilk hamle olarak P'yi (4,3)'e taşıyarak başlayabiliriz. Sen ne dersin?} \\
            }
        }
    }
     \\ \\

    \theutterance \stepcounter{utterance}  
    & & & \multicolumn{2}{p{0.3\linewidth}}{
        \cellcolor[rgb]{0.9,0.9,0.9}{
            \makecell[{{p{\linewidth}}}]{
                \texttt{\tiny{[GM$|$GM]}}
                \texttt{Öncelikle tahtamı kontrol edeyim. Önerdiğin hedef koordinatları boş mu diye bakıyorum:} \\
\\ 
\texttt{{-} P hedef (4,3): Benim tahtamda (4,3) boş ◌} \\
\texttt{{-} T hedef (5,4): Benim tahtamda (5,4) boş ◌} \\
\texttt{{-} S hedef (4,6): Benim tahtamda (4,6) boş ◌} \\
\texttt{{-} D hedef (3,4): Benim tahtamda (3,4) boş ◌} \\
\texttt{{-} A hedef (3,5): Benim tahtamda (3,5) boş ◌} \\
\texttt{{-} M hedef (4,4): Benim tahtamda (4,4) boş ◌} \\
\\ 
\texttt{Hepsi boş görünüyor. Şimdi O'lar için ne yapacağız? Onları olduğu yerde bırakıp diğerlerini bu hedeflere taşımaya başlayalım mı? Bence ilk hamle olarak P'yi (4,3)'e taşıyarak başlayabiliriz. Sen ne dersin?} \\
            }
        }
    }
    & & \\ \\

    \theutterance \stepcounter{utterance}  
    & & & \multicolumn{2}{p{0.3\linewidth}}{
        \cellcolor[rgb]{0.9,0.9,0.9}{
            \makecell[{{p{\linewidth}}}]{
                \texttt{\tiny{[GM$|$GM]}}
                \texttt{Invalid response format} \\
            }
        }
    }
    & & \\ \\

    \theutterance \stepcounter{utterance}  
    & & \multicolumn{4}{p{0.6\linewidth}}{
        \cellcolor[rgb]{0.9,0.9,0.9}{
            \makecell[{{p{\linewidth}}}]{
                \texttt{\tiny{[GM$\rangle$P2]}}
                \texttt{Ceza: Mesajın beklenen biçimde değil!} \\
\\ 
\texttt{Yanıtının yalnızca SÖYLE: <MESAJ> ya da HAMLE: <NESNE>, (<X>, <Y>) komutlarından birini içerdiğinden emin ol; bunun dışında hiçbir şey yazma!} \\
\\ 
\texttt{Toplamda 16 cezanın 3 kadarını biriktirdiniz.Lütfen tekrar deneyin!} \\
            }
        }
    }
    & \\ \\

    \theutterance \stepcounter{utterance}  
    & & & \multicolumn{4}{p{0.6\linewidth}}{
        \cellcolor[rgb]{0.9,0.9,0.9}{
            \makecell[{{p{\linewidth}}}]{
                \texttt{\tiny{[GM$\langle$P2]}}
                \texttt{SÖYLE: Hedef koordinatların hepsi benim tahtamda boş. O'ları olduğu yerde bırakalım. İlk hamle olarak P'yi (4,3)'e taşıyarak başlayalım. Sen de P'yi (4,3)'e taşıyabilir misin?} \\
            }
        }
    }
     \\ \\

    \theutterance \stepcounter{utterance}  
    & & & \multicolumn{2}{p{0.3\linewidth}}{
        \cellcolor[rgb]{0.9,0.9,0.9}{
            \makecell[{{p{\linewidth}}}]{
                \texttt{\tiny{[GM$|$GM]}}
                \texttt{SÖYLE: Hedef koordinatların hepsi benim tahtamda boş. O'ları olduğu yerde bırakalım. İlk hamle olarak P'yi (4,3)'e taşıyarak başlayalım. Sen de P'yi (4,3)'e taşıyabilir misin?} \\
            }
        }
    }
    & & \\ \\

    \theutterance \stepcounter{utterance}  
    & & \multicolumn{4}{p{0.6\linewidth}}{
        \cellcolor[rgb]{0.9,0.9,0.9}{
            \makecell[{{p{\linewidth}}}]{
                \texttt{\tiny{[P1$\langle$GM]}}
                \texttt{Mesajın diğer oyuncuya iletildi.} \\
\\ 
\\ 
\texttt{Şu anda en fazla 28 turdan 3. turu oynuyorsunuz.} \\
\texttt{Toplamda 16 cezanın 3 kadarını biriktirdiniz.Diğer oyuncu şu mesajı gönderdi:} \\
\\ 
\\ 
\texttt{"Hedef koordinatların hepsi benim tahtamda boş. O'ları olduğu yerde bırakalım. İlk hamle olarak P'yi (4,3)'e taşıyarak başlayalım. Sen de P'yi (4,3)'e taşıyabilir misin?"} \\
\\ 
\\ 
\texttt{Bir sonraki komutun nedir?} \\
            }
        }
    }
    & \\ \\

    \theutterance \stepcounter{utterance}  
    & \multicolumn{4}{p{0.6\linewidth}}{
        \cellcolor[rgb]{0.9,0.9,0.9}{
            \makecell[{{p{\linewidth}}}]{
                \texttt{\tiny{[P1$\rangle$GM]}}
                \texttt{HAMLE: P, (4,3)} \\
            }
        }
    }
    & & \\ \\

    \theutterance \stepcounter{utterance}  
    & & & \multicolumn{2}{p{0.3\linewidth}}{
        \cellcolor[rgb]{0.9,0.9,0.9}{
            \makecell[{{p{\linewidth}}}]{
                \texttt{\tiny{[GM$|$GM]}}
                \texttt{HAMLE: P, (4,3)} \\
            }
        }
    }
    & & \\ \\

    \theutterance \stepcounter{utterance}  
    & & & \multicolumn{2}{p{0.3\linewidth}}{
        \cellcolor[rgb]{0.9,0.9,0.9}{
            \makecell[{{p{\linewidth}}}]{
                \texttt{\tiny{[GM$|$GM]}}
                \texttt{Ceza: (4,3) boş değil; "O" içeriyor.} \\
\\ 
\texttt{Toplamda 16 cezanın 4 kadarını biriktirdiniz.Lütfen tekrar deneyin!} \\
            }
        }
    }
    & & \\ \\

    \theutterance \stepcounter{utterance}  
    & & \multicolumn{4}{p{0.6\linewidth}}{
        \cellcolor[rgb]{0.9,0.9,0.9}{
            \makecell[{{p{\linewidth}}}]{
                \texttt{\tiny{[P1$\langle$GM]}}
                \texttt{Ceza: (4,3) boş değil; "O" içeriyor.} \\
\\ 
\texttt{Toplamda 16 cezanın 4 kadarını biriktirdiniz.Lütfen tekrar deneyin!} \\
            }
        }
    }
    & \\ \\

    \theutterance \stepcounter{utterance}  
    & \multicolumn{4}{p{0.6\linewidth}}{
        \cellcolor[rgb]{0.9,0.9,0.9}{
            \makecell[{{p{\linewidth}}}]{
                \texttt{\tiny{[P1$\rangle$GM]}}
                \texttt{SÖYLE: (4,3) koordinatı benim tahtamda boş değil, O içeriyor. P için başka bir hedef seçelim. Benim tahtamda boş olan koordinatlar: (4,1), (5,1), (6,1), (7,1), (1,2), (2,2), (5,2), (7,2), (1,3), (2,3), (4,3) O var, (5,3), (6,3), (7,3), (1,4), (2,4), (5,4), (6,4), (7,4), (1,5), (4,5) P var, (5,5), (6,5), (7,5), (1,6), (2,6), (3,6), (4,6), (1,7), (2,7) S var, (3,7), (4,7), (5,7), (6,7), (7,7). P için (5,4) uygun mu?} \\
            }
        }
    }
    & & \\ \\

    \theutterance \stepcounter{utterance}  
    & & & \multicolumn{2}{p{0.3\linewidth}}{
        \cellcolor[rgb]{0.9,0.9,0.9}{
            \makecell[{{p{\linewidth}}}]{
                \texttt{\tiny{[GM$|$GM]}}
                \texttt{SÖYLE: (4,3) koordinatı benim tahtamda boş değil, O içeriyor. P için başka bir hedef seçelim. Benim tahtamda boş olan koordinatlar: (4,1), (5,1), (6,1), (7,1), (1,2), (2,2), (5,2), (7,2), (1,3), (2,3), (4,3) O var, (5,3), (6,3), (7,3), (1,4), (2,4), (5,4), (6,4), (7,4), (1,5), (4,5) P var, (5,5), (6,5), (7,5), (1,6), (2,6), (3,6), (4,6), (1,7), (2,7) S var, (3,7), (4,7), (5,7), (6,7), (7,7). P için (5,4) uygun mu?} \\
            }
        }
    }
    & & \\ \\

    \theutterance \stepcounter{utterance}  
    & & \multicolumn{4}{p{0.6\linewidth}}{
        \cellcolor[rgb]{0.9,0.9,0.9}{
            \makecell[{{p{\linewidth}}}]{
                \texttt{\tiny{[GM$\rangle$P2]}}
                \texttt{Mesajın diğer oyuncuya iletildi.} \\
\\ 
\\ 
\texttt{Şu anda en fazla 28 turdan 4. turu oynuyorsunuz.} \\
\texttt{Toplamda 16 cezanın 4 kadarını biriktirdiniz.Diğer oyuncu şu mesajı gönderdi:} \\
\\ 
\\ 
\texttt{"(4,3) koordinatı benim tahtamda boş değil, O içeriyor. P için başka bir hedef seçelim. Benim tahtamda boş olan koordinatlar: (4,1), (5,1), (6,1), (7,1), (1,2), (2,2), (5,2), (7,2), (1,3), (2,3), (4,3) O var, (5,3), (6,3), (7,3), (1,4), (2,4), (5,4), (6,4), (7,4), (1,5), (4,5) P var, (5,5), (6,5), (7,5), (1,6), (2,6), (3,6), (4,6), (1,7), (2,7) S var, (3,7), (4,7), (5,7), (6,7), (7,7). P için (5,4) uygun mu?"} \\
\\ 
\\ 
\texttt{Bir sonraki komutun nedir?} \\
            }
        }
    }
    & \\ \\

    \theutterance \stepcounter{utterance}  
    & & & \multicolumn{4}{p{0.6\linewidth}}{
        \cellcolor[rgb]{0.9,0.9,0.9}{
            \makecell[{{p{\linewidth}}}]{
                \texttt{\tiny{[GM$\langle$P2]}}
                \texttt{SÖYLE: Peki, P için (5,4) benim tahtamda boş, uygun. O zaman hedefleri güncelleyelim: P(5,4), T(?,?), S(?,?), D(?,?), A(?,?), M(?,?). Şimdi T için senin tahtanda boş olan bir koordinat öner. Benim tahtamda şu an tüm hedeflerin boş olması için dikkatli seçelim. Sana kendi tahtamdaki boş koordinatları da söyleyeyim mi?} \\
            }
        }
    }
     \\ \\

    \theutterance \stepcounter{utterance}  
    & & & \multicolumn{2}{p{0.3\linewidth}}{
        \cellcolor[rgb]{0.9,0.9,0.9}{
            \makecell[{{p{\linewidth}}}]{
                \texttt{\tiny{[GM$|$GM]}}
                \texttt{SÖYLE: Peki, P için (5,4) benim tahtamda boş, uygun. O zaman hedefleri güncelleyelim: P(5,4), T(?,?), S(?,?), D(?,?), A(?,?), M(?,?). Şimdi T için senin tahtanda boş olan bir koordinat öner. Benim tahtamda şu an tüm hedeflerin boş olması için dikkatli seçelim. Sana kendi tahtamdaki boş koordinatları da söyleyeyim mi?} \\
            }
        }
    }
    & & \\ \\

    \theutterance \stepcounter{utterance}  
    & & \multicolumn{4}{p{0.6\linewidth}}{
        \cellcolor[rgb]{0.9,0.9,0.9}{
            \makecell[{{p{\linewidth}}}]{
                \texttt{\tiny{[P1$\langle$GM]}}
                \texttt{Mesajın diğer oyuncuya iletildi.} \\
\\ 
\\ 
\texttt{Şu anda en fazla 28 turdan 4. turu oynuyorsunuz.} \\
\texttt{Toplamda 16 cezanın 4 kadarını biriktirdiniz.Diğer oyuncu şu mesajı gönderdi:} \\
\\ 
\\ 
\texttt{"Peki, P için (5,4) benim tahtamda boş, uygun. O zaman hedefleri güncelleyelim: P(5,4), T(?,?), S(?,?), D(?,?), A(?,?), M(?,?). Şimdi T için senin tahtanda boş olan bir koordinat öner. Benim tahtamda şu an tüm hedeflerin boş olması için dikkatli seçelim. Sana kendi tahtamdaki boş koordinatları da söyleyeyim mi?"} \\
\\ 
\\ 
\texttt{Bir sonraki komutun nedir?} \\
            }
        }
    }
    & \\ \\

    \theutterance \stepcounter{utterance}  
    & \multicolumn{4}{p{0.6\linewidth}}{
        \cellcolor[rgb]{0.9,0.9,0.9}{
            \makecell[{{p{\linewidth}}}]{
                \texttt{\tiny{[P1$\rangle$GM]}}
                \texttt{SÖYLE: Evet, lütfen tahtandaki boş koordinatları söyle. T için benim tahtamda şu koordinatlar boş: (1,1) zaten boş, (2,1) boş, (4,1), (5,1), (6,1), (7,1), (1,2), (2,2), (5,2), (7,2), (1,3), (2,3), (5,3), (6,3), (7,3), (1,4), (2,4), (6,4), (7,4), (1,5), (5,5), (6,5), (7,5), (1,6), (2,6), (3,6), (4,6), (1,7), (3,7), (4,7), (5,7), (6,7), (7,7). T için (6,4) öneriyorum, sende boş mu?} \\
            }
        }
    }
    & & \\ \\

    \theutterance \stepcounter{utterance}  
    & & & \multicolumn{2}{p{0.3\linewidth}}{
        \cellcolor[rgb]{0.9,0.9,0.9}{
            \makecell[{{p{\linewidth}}}]{
                \texttt{\tiny{[GM$|$GM]}}
                \texttt{SÖYLE: Evet, lütfen tahtandaki boş koordinatları söyle. T için benim tahtamda şu koordinatlar boş: (1,1) zaten boş, (2,1) boş, (4,1), (5,1), (6,1), (7,1), (1,2), (2,2), (5,2), (7,2), (1,3), (2,3), (5,3), (6,3), (7,3), (1,4), (2,4), (6,4), (7,4), (1,5), (5,5), (6,5), (7,5), (1,6), (2,6), (3,6), (4,6), (1,7), (3,7), (4,7), (5,7), (6,7), (7,7). T için (6,4) öneriyorum, sende boş mu?} \\
            }
        }
    }
    & & \\ \\

    \theutterance \stepcounter{utterance}  
    & & \multicolumn{4}{p{0.6\linewidth}}{
        \cellcolor[rgb]{0.9,0.9,0.9}{
            \makecell[{{p{\linewidth}}}]{
                \texttt{\tiny{[GM$\rangle$P2]}}
                \texttt{Mesajın diğer oyuncuya iletildi.} \\
\\ 
\\ 
\texttt{Şu anda en fazla 28 turdan 5. turu oynuyorsunuz.} \\
\texttt{Toplamda 16 cezanın 4 kadarını biriktirdiniz.Diğer oyuncu şu mesajı gönderdi:} \\
\\ 
\\ 
\texttt{"Evet, lütfen tahtandaki boş koordinatları söyle. T için benim tahtamda şu koordinatlar boş: (1,1) zaten boş, (2,1) boş, (4,1), (5,1), (6,1), (7,1), (1,2), (2,2), (5,2), (7,2), (1,3), (2,3), (5,3), (6,3), (7,3), (1,4), (2,4), (6,4), (7,4), (1,5), (5,5), (6,5), (7,5), (1,6), (2,6), (3,6), (4,6), (1,7), (3,7), (4,7), (5,7), (6,7), (7,7). T için (6,4) öneriyorum, sende boş mu?"} \\
\\ 
\\ 
\texttt{Bir sonraki komutun nedir?} \\
            }
        }
    }
    & \\ \\

    \theutterance \stepcounter{utterance}  
    & & & \multicolumn{4}{p{0.6\linewidth}}{
        \cellcolor[rgb]{0.9,0.9,0.9}{
            \makecell[{{p{\linewidth}}}]{
                \texttt{\tiny{[GM$\langle$P2]}}
                \texttt{SÖYLE: Benim tahtamda boş koordinatlar: (1,1), (2,1), (4,1), (5,1), (6,1), (7,1), (1,2) zaten O var, (2,2) boş, (3,2) O var, (5,2) boş, (7,2) boş, (1,3) boş, (2,3) boş, (4,3) boş, (6,3) boş, (7,3) boş, (1,4) boş, (2,4) boş, (3,4) boş, (4,4) boş, (5,4) boş, (6,4) boş, (7,4) boş, (1,5) boş, (2,5) boş, (4,5) boş, (5,5) boş, (6,5) boş, (7,5) boş, (1,6) boş, (2,6) boş, (3,6) boş, (4,6) boş, (5,6) boş, (7,6) boş, (1,7) M var, (2,7) boş, (3,7) boş, (4,7) boş, (5,7) boş, (6,7) S var, (7,7) T var. T için (6,4) bende boş, uygun. Hedefleri toparlayalım: P(5,4), T(6,4), diğerleri için sırayla gidelim. S için sen öner.} \\
            }
        }
    }
     \\ \\

    \theutterance \stepcounter{utterance}  
    & & & \multicolumn{2}{p{0.3\linewidth}}{
        \cellcolor[rgb]{0.9,0.9,0.9}{
            \makecell[{{p{\linewidth}}}]{
                \texttt{\tiny{[GM$|$GM]}}
                \texttt{SÖYLE: Benim tahtamda boş koordinatlar: (1,1), (2,1), (4,1), (5,1), (6,1), (7,1), (1,2) zaten O var, (2,2) boş, (3,2) O var, (5,2) boş, (7,2) boş, (1,3) boş, (2,3) boş, (4,3) boş, (6,3) boş, (7,3) boş, (1,4) boş, (2,4) boş, (3,4) boş, (4,4) boş, (5,4) boş, (6,4) boş, (7,4) boş, (1,5) boş, (2,5) boş, (4,5) boş, (5,5) boş, (6,5) boş, (7,5) boş, (1,6) boş, (2,6) boş, (3,6) boş, (4,6) boş, (5,6) boş, (7,6) boş, (1,7) M var, (2,7) boş, (3,7) boş, (4,7) boş, (5,7) boş, (6,7) S var, (7,7) T var. T için (6,4) bende boş, uygun. Hedefleri toparlayalım: P(5,4), T(6,4), diğerleri için sırayla gidelim. S için sen öner.} \\
            }
        }
    }
    & & \\ \\

    \theutterance \stepcounter{utterance}  
    & & \multicolumn{4}{p{0.6\linewidth}}{
        \cellcolor[rgb]{0.9,0.9,0.9}{
            \makecell[{{p{\linewidth}}}]{
                \texttt{\tiny{[P1$\langle$GM]}}
                \texttt{Mesajın diğer oyuncuya iletildi.} \\
\\ 
\\ 
\texttt{Şu anda en fazla 28 turdan 5. turu oynuyorsunuz.} \\
\texttt{Toplamda 16 cezanın 4 kadarını biriktirdiniz.Diğer oyuncu şu mesajı gönderdi:} \\
\\ 
\\ 
\texttt{"Benim tahtamda boş koordinatlar: (1,1), (2,1), (4,1), (5,1), (6,1), (7,1), (1,2) zaten O var, (2,2) boş, (3,2) O var, (5,2) boş, (7,2) boş, (1,3) boş, (2,3) boş, (4,3) boş, (6,3) boş, (7,3) boş, (1,4) boş, (2,4) boş, (3,4) boş, (4,4) boş, (5,4) boş, (6,4) boş, (7,4) boş, (1,5) boş, (2,5) boş, (4,5) boş, (5,5) boş, (6,5) boş, (7,5) boş, (1,6) boş, (2,6) boş, (3,6) boş, (4,6) boş, (5,6) boş, (7,6) boş, (1,7) M var, (2,7) boş, (3,7) boş, (4,7) boş, (5,7) boş, (6,7) S var, (7,7) T var. T için (6,4) bende boş, uygun. Hedefleri toparlayalım: P(5,4), T(6,4), diğerleri için sırayla gidelim. S için sen öner."} \\
\\ 
\\ 
\texttt{Bir sonraki komutun nedir?} \\
            }
        }
    }
    & \\ \\

    \theutterance \stepcounter{utterance}  
    & \multicolumn{4}{p{0.6\linewidth}}{
        \cellcolor[rgb]{0.9,0.9,0.9}{
            \makecell[{{p{\linewidth}}}]{
                \texttt{\tiny{[P1$\rangle$GM]}}
                \texttt{SÖYLE: S için benim tahtamda (4,6) boş, sende de boş görünüyor. S hedefi (4,6) olsun. D için (3,4) öneriyorum, sende boş mu?} \\
            }
        }
    }
    & & \\ \\

    \theutterance \stepcounter{utterance}  
    & & & \multicolumn{2}{p{0.3\linewidth}}{
        \cellcolor[rgb]{0.9,0.9,0.9}{
            \makecell[{{p{\linewidth}}}]{
                \texttt{\tiny{[GM$|$GM]}}
                \texttt{SÖYLE: S için benim tahtamda (4,6) boş, sende de boş görünüyor. S hedefi (4,6) olsun. D için (3,4) öneriyorum, sende boş mu?} \\
            }
        }
    }
    & & \\ \\

    \theutterance \stepcounter{utterance}  
    & & \multicolumn{4}{p{0.6\linewidth}}{
        \cellcolor[rgb]{0.9,0.9,0.9}{
            \makecell[{{p{\linewidth}}}]{
                \texttt{\tiny{[GM$\rangle$P2]}}
                \texttt{Mesajın diğer oyuncuya iletildi.} \\
\\ 
\\ 
\texttt{Şu anda en fazla 28 turdan 6. turu oynuyorsunuz.} \\
\texttt{Toplamda 16 cezanın 4 kadarını biriktirdiniz.Diğer oyuncu şu mesajı gönderdi:} \\
\\ 
\\ 
\texttt{"S için benim tahtamda (4,6) boş, sende de boş görünüyor. S hedefi (4,6) olsun. D için (3,4) öneriyorum, sende boş mu?"} \\
\\ 
\\ 
\texttt{Bir sonraki komutun nedir?} \\
            }
        }
    }
    & \\ \\

    \theutterance \stepcounter{utterance}  
    & & & \multicolumn{4}{p{0.6\linewidth}}{
        \cellcolor[rgb]{0.9,0.9,0.9}{
            \makecell[{{p{\linewidth}}}]{
                \texttt{\tiny{[GM$\langle$P2]}}
                \texttt{SÖYLE: S için (4,6) bende boş, uygun. D için (3,4) de bende boş. Hedefler: P(5,4), T(6,4), S(4,6), D(3,4). A için sen öner, M için de sırayla belirleyelim. Bu arada artık hamle yapmaya başlayalım mı? İlk hamle olarak P'yi (5,4)'e taşıyarak başlayabiliriz.} \\
            }
        }
    }
     \\ \\

    \theutterance \stepcounter{utterance}  
    & & & \multicolumn{2}{p{0.3\linewidth}}{
        \cellcolor[rgb]{0.9,0.9,0.9}{
            \makecell[{{p{\linewidth}}}]{
                \texttt{\tiny{[GM$|$GM]}}
                \texttt{SÖYLE: S için (4,6) bende boş, uygun. D için (3,4) de bende boş. Hedefler: P(5,4), T(6,4), S(4,6), D(3,4). A için sen öner, M için de sırayla belirleyelim. Bu arada artık hamle yapmaya başlayalım mı? İlk hamle olarak P'yi (5,4)'e taşıyarak başlayabiliriz.} \\
            }
        }
    }
    & & \\ \\

    \theutterance \stepcounter{utterance}  
    & & \multicolumn{4}{p{0.6\linewidth}}{
        \cellcolor[rgb]{0.9,0.9,0.9}{
            \makecell[{{p{\linewidth}}}]{
                \texttt{\tiny{[P1$\langle$GM]}}
                \texttt{Mesajın diğer oyuncuya iletildi.} \\
\\ 
\\ 
\texttt{Şu anda en fazla 28 turdan 6. turu oynuyorsunuz.} \\
\texttt{Toplamda 16 cezanın 4 kadarını biriktirdiniz.Diğer oyuncu şu mesajı gönderdi:} \\
\\ 
\\ 
\texttt{"S için (4,6) bende boş, uygun. D için (3,4) de bende boş. Hedefler: P(5,4), T(6,4), S(4,6), D(3,4). A için sen öner, M için de sırayla belirleyelim. Bu arada artık hamle yapmaya başlayalım mı? İlk hamle olarak P'yi (5,4)'e taşıyarak başlayabiliriz."} \\
\\ 
\\ 
\texttt{Bir sonraki komutun nedir?} \\
            }
        }
    }
    & \\ \\

    \theutterance \stepcounter{utterance}  
    & \multicolumn{4}{p{0.6\linewidth}}{
        \cellcolor[rgb]{0.9,0.9,0.9}{
            \makecell[{{p{\linewidth}}}]{
                \texttt{\tiny{[P1$\rangle$GM]}}
                \texttt{SÖYLE: A için (2,4) öneriyorum, sende boş mu? M için de (5,6) düşünüyorum. Hedefler: P(5,4), T(6,4), S(4,6), D(3,4), A(2,4), M(5,6). Hepsi sende boşsa P'yi (5,4)'e taşıyarak başlayalım.} \\
            }
        }
    }
    & & \\ \\

    \theutterance \stepcounter{utterance}  
    & & & \multicolumn{2}{p{0.3\linewidth}}{
        \cellcolor[rgb]{0.9,0.9,0.9}{
            \makecell[{{p{\linewidth}}}]{
                \texttt{\tiny{[GM$|$GM]}}
                \texttt{SÖYLE: A için (2,4) öneriyorum, sende boş mu? M için de (5,6) düşünüyorum. Hedefler: P(5,4), T(6,4), S(4,6), D(3,4), A(2,4), M(5,6). Hepsi sende boşsa P'yi (5,4)'e taşıyarak başlayalım.} \\
            }
        }
    }
    & & \\ \\

    \theutterance \stepcounter{utterance}  
    & & \multicolumn{4}{p{0.6\linewidth}}{
        \cellcolor[rgb]{0.9,0.9,0.9}{
            \makecell[{{p{\linewidth}}}]{
                \texttt{\tiny{[GM$\rangle$P2]}}
                \texttt{Mesajın diğer oyuncuya iletildi.} \\
\\ 
\\ 
\texttt{Şu anda en fazla 28 turdan 7. turu oynuyorsunuz.} \\
\texttt{Toplamda 16 cezanın 4 kadarını biriktirdiniz.Diğer oyuncu şu mesajı gönderdi:} \\
\\ 
\\ 
\texttt{"A için (2,4) öneriyorum, sende boş mu? M için de (5,6) düşünüyorum. Hedefler: P(5,4), T(6,4), S(4,6), D(3,4), A(2,4), M(5,6). Hepsi sende boşsa P'yi (5,4)'e taşıyarak başlayalım."} \\
\\ 
\\ 
\texttt{Bir sonraki komutun nedir?} \\
            }
        }
    }
    & \\ \\

    \theutterance \stepcounter{utterance}  
    & & & \multicolumn{4}{p{0.6\linewidth}}{
        \cellcolor[rgb]{0.9,0.9,0.9}{
            \makecell[{{p{\linewidth}}}]{
                \texttt{\tiny{[GM$\langle$P2]}}
                \texttt{SÖYLE: A için (2,4) bende boş, M için (5,6) da bende boş. Tüm hedefler onaylandı: P(5,4), T(6,4), S(4,6), D(3,4), A(2,4), M(5,6). Şimdi P'yi (5,4)'e taşıyorum. HAMLE: P, (5,4)} \\
            }
        }
    }
     \\ \\

    \theutterance \stepcounter{utterance}  
    & & & \multicolumn{2}{p{0.3\linewidth}}{
        \cellcolor[rgb]{0.9,0.9,0.9}{
            \makecell[{{p{\linewidth}}}]{
                \texttt{\tiny{[GM$|$GM]}}
                \texttt{SÖYLE: A için (2,4) bende boş, M için (5,6) da bende boş. Tüm hedefler onaylandı: P(5,4), T(6,4), S(4,6), D(3,4), A(2,4), M(5,6). Şimdi P'yi (5,4)'e taşıyorum. HAMLE: P, (5,4)} \\
            }
        }
    }
    & & \\ \\

    \theutterance \stepcounter{utterance}  
    & & & \multicolumn{2}{p{0.3\linewidth}}{
        \cellcolor[rgb]{0.9,0.9,0.9}{
            \makecell[{{p{\linewidth}}}]{
                \texttt{\tiny{[GM$|$GM]}}
                \texttt{Invalid response format: SÖYLE: A için (2,4) bende boş, M için (5,6) da bende boş. Tüm hedefler onaylandı: P(5,4), T(6,4), S(4,6), D(3,4), A(2,4), M(5,6). Şimdi P'yi (5,4)'e taşıyorum. HAMLE: P, (5,4)} \\
            }
        }
    }
    & & \\ \\

    \theutterance \stepcounter{utterance}  
    & & \multicolumn{4}{p{0.6\linewidth}}{
        \cellcolor[rgb]{0.9,0.9,0.9}{
            \makecell[{{p{\linewidth}}}]{
                \texttt{\tiny{[GM$\rangle$P2]}}
                \texttt{Ceza: Mesajında birden fazla komut var!} \\
\\ 
\texttt{Yanıtının yalnızca SÖYLE: <MESAJ> ya da HAMLE: <NESNE>, (<X>, <Y>) komutlarından birini içerdiğinden emin ol; bunun dışında hiçbir şey yazma!} \\
\\ 
\texttt{Toplamda 16 cezanın 5 kadarını biriktirdiniz.Lütfen tekrar deneyin!} \\
            }
        }
    }
    & \\ \\

    \theutterance \stepcounter{utterance}  
    & & & \multicolumn{4}{p{0.6\linewidth}}{
        \cellcolor[rgb]{0.9,0.9,0.9}{
            \makecell[{{p{\linewidth}}}]{
                \texttt{\tiny{[GM$\langle$P2]}}
                \texttt{HAMLE: P, (5,4)} \\
            }
        }
    }
     \\ \\

    \theutterance \stepcounter{utterance}  
    & & & \multicolumn{2}{p{0.3\linewidth}}{
        \cellcolor[rgb]{0.9,0.9,0.9}{
            \makecell[{{p{\linewidth}}}]{
                \texttt{\tiny{[GM$|$GM]}}
                \texttt{HAMLE: P, (5,4)} \\
            }
        }
    }
    & & \\ \\

    \theutterance \stepcounter{utterance}  
    & & & \multicolumn{2}{p{0.3\linewidth}}{
        \cellcolor[rgb]{0.9,0.9,0.9}{
            \makecell[{{p{\linewidth}}}]{
                \texttt{\tiny{[GM$|$GM]}}
                \texttt{P moved to (5, 4)} \\
            }
        }
    }
    & & \\ \\

    \theutterance \stepcounter{utterance}  
    & & \multicolumn{4}{p{0.6\linewidth}}{
        \cellcolor[rgb]{0.9,0.9,0.9}{
            \makecell[{{p{\linewidth}}}]{
                \texttt{\tiny{[P1$\langle$GM]}}
                \texttt{Mesajın diğer oyuncuya iletildi.} \\
\\ 
\\ 
\texttt{Şu anda en fazla 28 turdan 7. turu oynuyorsunuz.} \\
\texttt{Toplamda 16 cezanın 5 kadarını biriktirdiniz.Diğer oyuncu kendi ızgarasında bir nesneyi hareket ettirdi.} \\
\\ 
\\ 
\texttt{Bir sonraki komutun nedir?} \\
            }
        }
    }
    & \\ \\

    \theutterance \stepcounter{utterance}  
    & \multicolumn{4}{p{0.6\linewidth}}{
        \cellcolor[rgb]{0.9,0.9,0.9}{
            \makecell[{{p{\linewidth}}}]{
                \texttt{\tiny{[P1$\rangle$GM]}}
                \texttt{HAMLE: P, (5,4)} \\
            }
        }
    }
    & & \\ \\

    \theutterance \stepcounter{utterance}  
    & & & \multicolumn{2}{p{0.3\linewidth}}{
        \cellcolor[rgb]{0.9,0.9,0.9}{
            \makecell[{{p{\linewidth}}}]{
                \texttt{\tiny{[GM$|$GM]}}
                \texttt{HAMLE: P, (5,4)} \\
            }
        }
    }
    & & \\ \\

    \theutterance \stepcounter{utterance}  
    & & & \multicolumn{2}{p{0.3\linewidth}}{
        \cellcolor[rgb]{0.9,0.9,0.9}{
            \makecell[{{p{\linewidth}}}]{
                \texttt{\tiny{[GM$|$GM]}}
                \texttt{P moved to (5, 4)} \\
            }
        }
    }
    & & \\ \\

    \theutterance \stepcounter{utterance}  
    & & \multicolumn{4}{p{0.6\linewidth}}{
        \cellcolor[rgb]{0.9,0.9,0.9}{
            \makecell[{{p{\linewidth}}}]{
                \texttt{\tiny{[GM$\rangle$P2]}}
                \texttt{P’ başarıyla (5,4) koordinatlarına taşındı. Güncellenmiş ızgaran şu şekilde:} \\
\texttt{```} \\
\texttt{ 1234567} \\
\texttt{╔═╤══╤══╗} \\
\texttt{║◌│O◌│◌◌║ 1} \\
\texttt{║◌└┬─┴──╢ 2} \\
\texttt{║◌◌│◌D◌◌║ 3} \\
\texttt{╟─┬┘◌P┌─╢ 4} \\
\texttt{║◌│◌◌◌│A║ 5} \\
\texttt{╟─┼─┬─┤S║ 6} \\
\texttt{║M│◌│◌│T║ 7} \\
\texttt{╚═╧═╧═╧═╝} \\
\texttt{```} \\
\texttt{Şu anda en fazla 28 turdan 8. turu oynuyorsunuz.} \\
\texttt{Toplamda 16 cezanın 5 kadarını biriktirdiniz.Diğer oyuncu kendi ızgarasında bir nesneyi hareket ettirdi.} \\
\\ 
\\ 
\texttt{Bir sonraki komutun nedir?} \\
            }
        }
    }
    & \\ \\

    \theutterance \stepcounter{utterance}  
    & & & \multicolumn{4}{p{0.6\linewidth}}{
        \cellcolor[rgb]{0.9,0.9,0.9}{
            \makecell[{{p{\linewidth}}}]{
                \texttt{\tiny{[GM$\langle$P2]}}
                \texttt{SÖYLE: P'yi (5,4)'e taşıdım. Şimdi sırada T var, T'yi (6,4)'e taşıyorum.} \\
            }
        }
    }
     \\ \\

    \theutterance \stepcounter{utterance}  
    & & & \multicolumn{2}{p{0.3\linewidth}}{
        \cellcolor[rgb]{0.9,0.9,0.9}{
            \makecell[{{p{\linewidth}}}]{
                \texttt{\tiny{[GM$|$GM]}}
                \texttt{SÖYLE: P'yi (5,4)'e taşıdım. Şimdi sırada T var, T'yi (6,4)'e taşıyorum.} \\
            }
        }
    }
    & & \\ \\

    \theutterance \stepcounter{utterance}  
    & & \multicolumn{4}{p{0.6\linewidth}}{
        \cellcolor[rgb]{0.9,0.9,0.9}{
            \makecell[{{p{\linewidth}}}]{
                \texttt{\tiny{[P1$\langle$GM]}}
                \texttt{P’ başarıyla (5,4) koordinatlarına taşındı. Güncellenmiş ızgaran şu şekilde:} \\
\texttt{```} \\
\texttt{ 1234567} \\
\texttt{╔═╤══╤══╗} \\
\texttt{║◌│◌T│◌◌║ 1} \\
\texttt{║◌└┬─┴──╢ 2} \\
\texttt{║◌◌│O◌◌◌║ 3} \\
\texttt{╟─┬┘◌P┌─╢ 4} \\
\texttt{║◌│AD◌│◌║ 5} \\
\texttt{╟─┼─┬─┤M║ 6} \\
\texttt{║◌│S│◌│◌║ 7} \\
\texttt{╚═╧═╧═╧═╝} \\
\texttt{```} \\
\texttt{Şu anda en fazla 28 turdan 8. turu oynuyorsunuz.} \\
\texttt{Toplamda 16 cezanın 5 kadarını biriktirdiniz.Diğer oyuncu şu mesajı gönderdi:} \\
\\ 
\\ 
\texttt{"P'yi (5,4)'e taşıdım. Şimdi sırada T var, T'yi (6,4)'e taşıyorum."} \\
\\ 
\\ 
\texttt{Bir sonraki komutun nedir?} \\
            }
        }
    }
    & \\ \\

    \theutterance \stepcounter{utterance}  
    & \multicolumn{4}{p{0.6\linewidth}}{
        \cellcolor[rgb]{0.9,0.9,0.9}{
            \makecell[{{p{\linewidth}}}]{
                \texttt{\tiny{[P1$\rangle$GM]}}
                \texttt{SÖYLE: Tamam, ben de T'yi (6,4)'e taşıyorum. Sonra sırayla D, A, S, M şeklinde devam edelim.} \\
            }
        }
    }
    & & \\ \\

    \theutterance \stepcounter{utterance}  
    & & & \multicolumn{2}{p{0.3\linewidth}}{
        \cellcolor[rgb]{0.9,0.9,0.9}{
            \makecell[{{p{\linewidth}}}]{
                \texttt{\tiny{[GM$|$GM]}}
                \texttt{SÖYLE: Tamam, ben de T'yi (6,4)'e taşıyorum. Sonra sırayla D, A, S, M şeklinde devam edelim.} \\
            }
        }
    }
    & & \\ \\

    \theutterance \stepcounter{utterance}  
    & & \multicolumn{4}{p{0.6\linewidth}}{
        \cellcolor[rgb]{0.9,0.9,0.9}{
            \makecell[{{p{\linewidth}}}]{
                \texttt{\tiny{[GM$\rangle$P2]}}
                \texttt{Mesajın diğer oyuncuya iletildi.} \\
\\ 
\\ 
\texttt{Şu anda en fazla 28 turdan 9. turu oynuyorsunuz.} \\
\texttt{Toplamda 16 cezanın 5 kadarını biriktirdiniz.Diğer oyuncu şu mesajı gönderdi:} \\
\\ 
\\ 
\texttt{"Tamam, ben de T'yi (6,4)'e taşıyorum. Sonra sırayla D, A, S, M şeklinde devam edelim."} \\
\\ 
\\ 
\texttt{Bir sonraki komutun nedir?} \\
            }
        }
    }
    & \\ \\

    \theutterance \stepcounter{utterance}  
    & & & \multicolumn{4}{p{0.6\linewidth}}{
        \cellcolor[rgb]{0.9,0.9,0.9}{
            \makecell[{{p{\linewidth}}}]{
                \texttt{\tiny{[GM$\langle$P2]}}
                \texttt{HAMLE: T, (6,4)} \\
            }
        }
    }
     \\ \\

    \theutterance \stepcounter{utterance}  
    & & & \multicolumn{2}{p{0.3\linewidth}}{
        \cellcolor[rgb]{0.9,0.9,0.9}{
            \makecell[{{p{\linewidth}}}]{
                \texttt{\tiny{[GM$|$GM]}}
                \texttt{HAMLE: T, (6,4)} \\
            }
        }
    }
    & & \\ \\

    \theutterance \stepcounter{utterance}  
    & & & \multicolumn{2}{p{0.3\linewidth}}{
        \cellcolor[rgb]{0.9,0.9,0.9}{
            \makecell[{{p{\linewidth}}}]{
                \texttt{\tiny{[GM$|$GM]}}
                \texttt{Ceza: (6,4) boş değil; "┌" içeriyor.} \\
\\ 
\texttt{Toplamda 16 cezanın 6 kadarını biriktirdiniz.Lütfen tekrar deneyin!} \\
            }
        }
    }
    & & \\ \\

    \theutterance \stepcounter{utterance}  
    & & \multicolumn{4}{p{0.6\linewidth}}{
        \cellcolor[rgb]{0.9,0.9,0.9}{
            \makecell[{{p{\linewidth}}}]{
                \texttt{\tiny{[GM$\rangle$P2]}}
                \texttt{Ceza: (6,4) boş değil; "┌" içeriyor.} \\
\\ 
\texttt{Toplamda 16 cezanın 6 kadarını biriktirdiniz.Lütfen tekrar deneyin!} \\
            }
        }
    }
    & \\ \\

    \theutterance \stepcounter{utterance}  
    & & & \multicolumn{4}{p{0.6\linewidth}}{
        \cellcolor[rgb]{0.9,0.9,0.9}{
            \makecell[{{p{\linewidth}}}]{
                \texttt{\tiny{[GM$\langle$P2]}}
                \texttt{Tahtamı kontrol edeyim. (6,4) koordinatında ┌ sembolü var, bu çerçeve/çizgi karakteri. Boş değilmiş. Hedef olarak (6,4) yerine boş bir koordinat bulmalıyız. T için alternatif olarak (7,4) boş mu diye bakıyorum, evet bende boş. Ama önce partnerime sormalıyım çünkü onun tahtasında da boş olması lazım.} \\
\\ 
\texttt{SÖYLE: (6,4) benim tahtamda boş değil, ┌ sembolü var. T için alternatif hedef olarak (7,4) boş, sende de boş mu?} \\
            }
        }
    }
     \\ \\

    \theutterance \stepcounter{utterance}  
    & & & \multicolumn{2}{p{0.3\linewidth}}{
        \cellcolor[rgb]{0.9,0.9,0.9}{
            \makecell[{{p{\linewidth}}}]{
                \texttt{\tiny{[GM$|$GM]}}
                \texttt{Tahtamı kontrol edeyim. (6,4) koordinatında ┌ sembolü var, bu çerçeve/çizgi karakteri. Boş değilmiş. Hedef olarak (6,4) yerine boş bir koordinat bulmalıyız. T için alternatif olarak (7,4) boş mu diye bakıyorum, evet bende boş. Ama önce partnerime sormalıyım çünkü onun tahtasında da boş olması lazım.} \\
\\ 
\texttt{SÖYLE: (6,4) benim tahtamda boş değil, ┌ sembolü var. T için alternatif hedef olarak (7,4) boş, sende de boş mu?} \\
            }
        }
    }
    & & \\ \\

    \theutterance \stepcounter{utterance}  
    & & & \multicolumn{2}{p{0.3\linewidth}}{
        \cellcolor[rgb]{0.9,0.9,0.9}{
            \makecell[{{p{\linewidth}}}]{
                \texttt{\tiny{[GM$|$GM]}}
                \texttt{Invalid format: head is not empty: 'Tahtamı kontrol edeyim. (6,4) koordinatında ┌ sembolü var, bu çerçeve/çizgi karakteri. Boş değilmiş. Hedef olarak (6,4) yerine boş bir koordinat bulmalıyız. T için alternatif olarak (7,4) boş mu diye bakıyorum, evet bende boş. Ama önce partnerime sormalıyım çünkü onun tahtasında da boş olması lazım.} \\
\\ 
\texttt{'} \\
            }
        }
    }
    & & \\ \\

    \theutterance \stepcounter{utterance}  
    & & \multicolumn{4}{p{0.6\linewidth}}{
        \cellcolor[rgb]{0.9,0.9,0.9}{
            \makecell[{{p{\linewidth}}}]{
                \texttt{\tiny{[GM$\rangle$P2]}}
                \texttt{Ceza: Mesajında komuttan önce hiçbir şey bulunmamalı!} \\
\\ 
\texttt{Yanıtının yalnızca SÖYLE: <MESAJ> ya da HAMLE: <NESNE>, (<X>, <Y>) komutlarından birini içerdiğinden emin ol; bunun dışında hiçbir şey yazma!} \\
\\ 
\texttt{Toplamda 16 cezanın 7 kadarını biriktirdiniz.Lütfen tekrar deneyin!} \\
            }
        }
    }
    & \\ \\

    \theutterance \stepcounter{utterance}  
    & & & \multicolumn{4}{p{0.6\linewidth}}{
        \cellcolor[rgb]{0.9,0.9,0.9}{
            \makecell[{{p{\linewidth}}}]{
                \texttt{\tiny{[GM$\langle$P2]}}
                \texttt{SÖYLE: (6,4) benim tahtamda boş değil, ┌ sembolü var. T için alternatif hedef olarak (7,4) boş, sende de boş mu?} \\
            }
        }
    }
     \\ \\

    \theutterance \stepcounter{utterance}  
    & & & \multicolumn{2}{p{0.3\linewidth}}{
        \cellcolor[rgb]{0.9,0.9,0.9}{
            \makecell[{{p{\linewidth}}}]{
                \texttt{\tiny{[GM$|$GM]}}
                \texttt{SÖYLE: (6,4) benim tahtamda boş değil, ┌ sembolü var. T için alternatif hedef olarak (7,4) boş, sende de boş mu?} \\
            }
        }
    }
    & & \\ \\

    \theutterance \stepcounter{utterance}  
    & & \multicolumn{4}{p{0.6\linewidth}}{
        \cellcolor[rgb]{0.9,0.9,0.9}{
            \makecell[{{p{\linewidth}}}]{
                \texttt{\tiny{[P1$\langle$GM]}}
                \texttt{Mesajın diğer oyuncuya iletildi.} \\
\\ 
\\ 
\texttt{Şu anda en fazla 28 turdan 9. turu oynuyorsunuz.} \\
\texttt{Toplamda 16 cezanın 7 kadarını biriktirdiniz.Diğer oyuncu şu mesajı gönderdi:} \\
\\ 
\\ 
\texttt{"(6,4) benim tahtamda boş değil, ┌ sembolü var. T için alternatif hedef olarak (7,4) boş, sende de boş mu?"} \\
\\ 
\\ 
\texttt{Bir sonraki komutun nedir?} \\
            }
        }
    }
    & \\ \\

    \theutterance \stepcounter{utterance}  
    & \multicolumn{4}{p{0.6\linewidth}}{
        \cellcolor[rgb]{0.9,0.9,0.9}{
            \makecell[{{p{\linewidth}}}]{
                \texttt{\tiny{[P1$\rangle$GM]}}
                \texttt{SÖYLE: (7,4) benim tahtamda boş, uygun. T hedefini (7,4) olarak güncelleyelim. O zaman yeni hedefler: P(5,4) tamam, T(7,4), S(4,6), D(3,4), A(2,4), M(5,6). T'yi (7,4)'e taşıyorum.} \\
            }
        }
    }
    & & \\ \\

    \theutterance \stepcounter{utterance}  
    & & & \multicolumn{2}{p{0.3\linewidth}}{
        \cellcolor[rgb]{0.9,0.9,0.9}{
            \makecell[{{p{\linewidth}}}]{
                \texttt{\tiny{[GM$|$GM]}}
                \texttt{SÖYLE: (7,4) benim tahtamda boş, uygun. T hedefini (7,4) olarak güncelleyelim. O zaman yeni hedefler: P(5,4) tamam, T(7,4), S(4,6), D(3,4), A(2,4), M(5,6). T'yi (7,4)'e taşıyorum.} \\
            }
        }
    }
    & & \\ \\

    \theutterance \stepcounter{utterance}  
    & & \multicolumn{4}{p{0.6\linewidth}}{
        \cellcolor[rgb]{0.9,0.9,0.9}{
            \makecell[{{p{\linewidth}}}]{
                \texttt{\tiny{[GM$\rangle$P2]}}
                \texttt{Mesajın diğer oyuncuya iletildi.} \\
\\ 
\\ 
\texttt{Şu anda en fazla 28 turdan 10. turu oynuyorsunuz.} \\
\texttt{Toplamda 16 cezanın 7 kadarını biriktirdiniz.Diğer oyuncu şu mesajı gönderdi:} \\
\\ 
\\ 
\texttt{"(7,4) benim tahtamda boş, uygun. T hedefini (7,4) olarak güncelleyelim. O zaman yeni hedefler: P(5,4) tamam, T(7,4), S(4,6), D(3,4), A(2,4), M(5,6). T'yi (7,4)'e taşıyorum."} \\
\\ 
\\ 
\texttt{Bir sonraki komutun nedir?} \\
            }
        }
    }
    & \\ \\

    \theutterance \stepcounter{utterance}  
    & & & \multicolumn{4}{p{0.6\linewidth}}{
        \cellcolor[rgb]{0.9,0.9,0.9}{
            \makecell[{{p{\linewidth}}}]{
                \texttt{\tiny{[GM$\langle$P2]}}
                \texttt{HAMLE: T, (7,4)} \\
            }
        }
    }
     \\ \\

    \theutterance \stepcounter{utterance}  
    & & & \multicolumn{2}{p{0.3\linewidth}}{
        \cellcolor[rgb]{0.9,0.9,0.9}{
            \makecell[{{p{\linewidth}}}]{
                \texttt{\tiny{[GM$|$GM]}}
                \texttt{HAMLE: T, (7,4)} \\
            }
        }
    }
    & & \\ \\

    \theutterance \stepcounter{utterance}  
    & & & \multicolumn{2}{p{0.3\linewidth}}{
        \cellcolor[rgb]{0.9,0.9,0.9}{
            \makecell[{{p{\linewidth}}}]{
                \texttt{\tiny{[GM$|$GM]}}
                \texttt{Ceza: (7,4) boş değil; "─" içeriyor.} \\
\\ 
\texttt{Toplamda 16 cezanın 8 kadarını biriktirdiniz.Lütfen tekrar deneyin!} \\
            }
        }
    }
    & & \\ \\

    \theutterance \stepcounter{utterance}  
    & & \multicolumn{4}{p{0.6\linewidth}}{
        \cellcolor[rgb]{0.9,0.9,0.9}{
            \makecell[{{p{\linewidth}}}]{
                \texttt{\tiny{[GM$\rangle$P2]}}
                \texttt{Ceza: (7,4) boş değil; "─" içeriyor.} \\
\\ 
\texttt{Toplamda 16 cezanın 8 kadarını biriktirdiniz.Lütfen tekrar deneyin!} \\
            }
        }
    }
    & \\ \\

    \theutterance \stepcounter{utterance}  
    & & & \multicolumn{4}{p{0.6\linewidth}}{
        \cellcolor[rgb]{0.9,0.9,0.9}{
            \makecell[{{p{\linewidth}}}]{
                \texttt{\tiny{[GM$\langle$P2]}}
                \texttt{Tahtamı dikkatlice kontrol etmem gerekiyor. Izgarayı satır satır inceleyeyim.} \\
\\ 
\texttt{SÖYLE: (7,4) de boş değil, ─ sembolü var. T için boş koordinat bulmakta zorlanıyorum. Benim tahtamda boş olan ve senin de boş dediğin bir koordinat bulalım. (7,3) boş mu? Veya (6,3)? Lütfen kontrol et.} \\
            }
        }
    }
     \\ \\

    \theutterance \stepcounter{utterance}  
    & & & \multicolumn{2}{p{0.3\linewidth}}{
        \cellcolor[rgb]{0.9,0.9,0.9}{
            \makecell[{{p{\linewidth}}}]{
                \texttt{\tiny{[GM$|$GM]}}
                \texttt{Tahtamı dikkatlice kontrol etmem gerekiyor. Izgarayı satır satır inceleyeyim.} \\
\\ 
\texttt{SÖYLE: (7,4) de boş değil, ─ sembolü var. T için boş koordinat bulmakta zorlanıyorum. Benim tahtamda boş olan ve senin de boş dediğin bir koordinat bulalım. (7,3) boş mu? Veya (6,3)? Lütfen kontrol et.} \\
            }
        }
    }
    & & \\ \\

    \theutterance \stepcounter{utterance}  
    & & & \multicolumn{2}{p{0.3\linewidth}}{
        \cellcolor[rgb]{0.9,0.9,0.9}{
            \makecell[{{p{\linewidth}}}]{
                \texttt{\tiny{[GM$|$GM]}}
                \texttt{Invalid format: head is not empty: 'Tahtamı dikkatlice kontrol etmem gerekiyor. Izgarayı satır satır inceleyeyim.} \\
\\ 
\texttt{'} \\
            }
        }
    }
    & & \\ \\

    \theutterance \stepcounter{utterance}  
    & & \multicolumn{4}{p{0.6\linewidth}}{
        \cellcolor[rgb]{0.9,0.9,0.9}{
            \makecell[{{p{\linewidth}}}]{
                \texttt{\tiny{[GM$\rangle$P2]}}
                \texttt{Ceza: Mesajında komuttan önce hiçbir şey bulunmamalı!} \\
\\ 
\texttt{Yanıtının yalnızca SÖYLE: <MESAJ> ya da HAMLE: <NESNE>, (<X>, <Y>) komutlarından birini içerdiğinden emin ol; bunun dışında hiçbir şey yazma!} \\
\\ 
\texttt{Toplamda 16 cezanın 9 kadarını biriktirdiniz.Lütfen tekrar deneyin!} \\
            }
        }
    }
    & \\ \\

    \theutterance \stepcounter{utterance}  
    & & & \multicolumn{4}{p{0.6\linewidth}}{
        \cellcolor[rgb]{0.9,0.9,0.9}{
            \makecell[{{p{\linewidth}}}]{
                \texttt{\tiny{[GM$\langle$P2]}}
                \texttt{SÖYLE: (7,4) de boş değil, ─ sembolü var. T için boş koordinat bulmalıyız. Benim tahtamda (7,3) veya (6,3) boş görünüyor. Sende bunlar boş mu?} \\
            }
        }
    }
     \\ \\

    \theutterance \stepcounter{utterance}  
    & & & \multicolumn{2}{p{0.3\linewidth}}{
        \cellcolor[rgb]{0.9,0.9,0.9}{
            \makecell[{{p{\linewidth}}}]{
                \texttt{\tiny{[GM$|$GM]}}
                \texttt{SÖYLE: (7,4) de boş değil, ─ sembolü var. T için boş koordinat bulmalıyız. Benim tahtamda (7,3) veya (6,3) boş görünüyor. Sende bunlar boş mu?} \\
            }
        }
    }
    & & \\ \\

    \theutterance \stepcounter{utterance}  
    & & \multicolumn{4}{p{0.6\linewidth}}{
        \cellcolor[rgb]{0.9,0.9,0.9}{
            \makecell[{{p{\linewidth}}}]{
                \texttt{\tiny{[P1$\langle$GM]}}
                \texttt{Mesajın diğer oyuncuya iletildi.} \\
\\ 
\\ 
\texttt{Şu anda en fazla 28 turdan 10. turu oynuyorsunuz.} \\
\texttt{Toplamda 16 cezanın 9 kadarını biriktirdiniz.Diğer oyuncu şu mesajı gönderdi:} \\
\\ 
\\ 
\texttt{"(7,4) de boş değil, ─ sembolü var. T için boş koordinat bulmalıyız. Benim tahtamda (7,3) veya (6,3) boş görünüyor. Sende bunlar boş mu?"} \\
\\ 
\\ 
\texttt{Bir sonraki komutun nedir?} \\
            }
        }
    }
    & \\ \\

    \theutterance \stepcounter{utterance}  
    & \multicolumn{4}{p{0.6\linewidth}}{
        \cellcolor[rgb]{0.9,0.9,0.9}{
            \makecell[{{p{\linewidth}}}]{
                \texttt{\tiny{[P1$\rangle$GM]}}
                \texttt{SÖYLE: (7,3) benim tahtamda boş. T hedefini (7,3) yapalım. Güncel hedefler: P(5,4), T(7,3), S(4,6), D(3,4), A(2,4), M(5,6). T'yi (7,3)'e taşıyorum şimdi.} \\
            }
        }
    }
    & & \\ \\

    \theutterance \stepcounter{utterance}  
    & & & \multicolumn{2}{p{0.3\linewidth}}{
        \cellcolor[rgb]{0.9,0.9,0.9}{
            \makecell[{{p{\linewidth}}}]{
                \texttt{\tiny{[GM$|$GM]}}
                \texttt{SÖYLE: (7,3) benim tahtamda boş. T hedefini (7,3) yapalım. Güncel hedefler: P(5,4), T(7,3), S(4,6), D(3,4), A(2,4), M(5,6). T'yi (7,3)'e taşıyorum şimdi.} \\
            }
        }
    }
    & & \\ \\

    \theutterance \stepcounter{utterance}  
    & & \multicolumn{4}{p{0.6\linewidth}}{
        \cellcolor[rgb]{0.9,0.9,0.9}{
            \makecell[{{p{\linewidth}}}]{
                \texttt{\tiny{[GM$\rangle$P2]}}
                \texttt{Mesajın diğer oyuncuya iletildi.} \\
\\ 
\\ 
\texttt{Şu anda en fazla 28 turdan 11. turu oynuyorsunuz.} \\
\texttt{Toplamda 16 cezanın 9 kadarını biriktirdiniz.Diğer oyuncu şu mesajı gönderdi:} \\
\\ 
\\ 
\texttt{"(7,3) benim tahtamda boş. T hedefini (7,3) yapalım. Güncel hedefler: P(5,4), T(7,3), S(4,6), D(3,4), A(2,4), M(5,6). T'yi (7,3)'e taşıyorum şimdi."} \\
\\ 
\\ 
\texttt{Bir sonraki komutun nedir?} \\
            }
        }
    }
    & \\ \\

    \theutterance \stepcounter{utterance}  
    & & & \multicolumn{4}{p{0.6\linewidth}}{
        \cellcolor[rgb]{0.9,0.9,0.9}{
            \makecell[{{p{\linewidth}}}]{
                \texttt{\tiny{[GM$\langle$P2]}}
                \texttt{HAMLE: T, (7,3)} \\
            }
        }
    }
     \\ \\

    \theutterance \stepcounter{utterance}  
    & & & \multicolumn{2}{p{0.3\linewidth}}{
        \cellcolor[rgb]{0.9,0.9,0.9}{
            \makecell[{{p{\linewidth}}}]{
                \texttt{\tiny{[GM$|$GM]}}
                \texttt{HAMLE: T, (7,3)} \\
            }
        }
    }
    & & \\ \\

    \theutterance \stepcounter{utterance}  
    & & & \multicolumn{2}{p{0.3\linewidth}}{
        \cellcolor[rgb]{0.9,0.9,0.9}{
            \makecell[{{p{\linewidth}}}]{
                \texttt{\tiny{[GM$|$GM]}}
                \texttt{T moved to (7, 3)} \\
            }
        }
    }
    & & \\ \\

    \theutterance \stepcounter{utterance}  
    & & \multicolumn{4}{p{0.6\linewidth}}{
        \cellcolor[rgb]{0.9,0.9,0.9}{
            \makecell[{{p{\linewidth}}}]{
                \texttt{\tiny{[P1$\langle$GM]}}
                \texttt{Mesajın diğer oyuncuya iletildi.} \\
\\ 
\\ 
\texttt{Şu anda en fazla 28 turdan 11. turu oynuyorsunuz.} \\
\texttt{Toplamda 16 cezanın 9 kadarını biriktirdiniz.Diğer oyuncu kendi ızgarasında bir nesneyi hareket ettirdi.} \\
\\ 
\\ 
\texttt{Bir sonraki komutun nedir?} \\
            }
        }
    }
    & \\ \\

    \theutterance \stepcounter{utterance}  
    & \multicolumn{4}{p{0.6\linewidth}}{
        \cellcolor[rgb]{0.9,0.9,0.9}{
            \makecell[{{p{\linewidth}}}]{
                \texttt{\tiny{[P1$\rangle$GM]}}
                \texttt{HAMLE: T, (7,3)} \\
            }
        }
    }
    & & \\ \\

    \theutterance \stepcounter{utterance}  
    & & & \multicolumn{2}{p{0.3\linewidth}}{
        \cellcolor[rgb]{0.9,0.9,0.9}{
            \makecell[{{p{\linewidth}}}]{
                \texttt{\tiny{[GM$|$GM]}}
                \texttt{HAMLE: T, (7,3)} \\
            }
        }
    }
    & & \\ \\

    \theutterance \stepcounter{utterance}  
    & & & \multicolumn{2}{p{0.3\linewidth}}{
        \cellcolor[rgb]{0.9,0.9,0.9}{
            \makecell[{{p{\linewidth}}}]{
                \texttt{\tiny{[GM$|$GM]}}
                \texttt{T moved to (7, 3)} \\
            }
        }
    }
    & & \\ \\

    \theutterance \stepcounter{utterance}  
    & & \multicolumn{4}{p{0.6\linewidth}}{
        \cellcolor[rgb]{0.9,0.9,0.9}{
            \makecell[{{p{\linewidth}}}]{
                \texttt{\tiny{[GM$\rangle$P2]}}
                \texttt{T’ başarıyla (7,3) koordinatlarına taşındı. Güncellenmiş ızgaran şu şekilde:} \\
\texttt{```} \\
\texttt{ 1234567} \\
\texttt{╔═╤══╤══╗} \\
\texttt{║◌│O◌│◌◌║ 1} \\
\texttt{║◌└┬─┴──╢ 2} \\
\texttt{║◌◌│◌D◌T║ 3} \\
\texttt{╟─┬┘◌P┌─╢ 4} \\
\texttt{║◌│◌◌◌│A║ 5} \\
\texttt{╟─┼─┬─┤S║ 6} \\
\texttt{║M│◌│◌│◌║ 7} \\
\texttt{╚═╧═╧═╧═╝} \\
\texttt{```} \\
\texttt{Şu anda en fazla 28 turdan 12. turu oynuyorsunuz.} \\
\texttt{Toplamda 16 cezanın 9 kadarını biriktirdiniz.Diğer oyuncu kendi ızgarasında bir nesneyi hareket ettirdi.} \\
\\ 
\\ 
\texttt{Bir sonraki komutun nedir?} \\
            }
        }
    }
    & \\ \\

    \theutterance \stepcounter{utterance}  
    & & & \multicolumn{4}{p{0.6\linewidth}}{
        \cellcolor[rgb]{0.9,0.9,0.9}{
            \makecell[{{p{\linewidth}}}]{
                \texttt{\tiny{[GM$\langle$P2]}}
                \texttt{SÖYLE: T'yi (7,3)'e taşıdım. Şimdi sırada D var, D'yi (3,4)'e taşıyorum.} \\
            }
        }
    }
     \\ \\

    \theutterance \stepcounter{utterance}  
    & & & \multicolumn{2}{p{0.3\linewidth}}{
        \cellcolor[rgb]{0.9,0.9,0.9}{
            \makecell[{{p{\linewidth}}}]{
                \texttt{\tiny{[GM$|$GM]}}
                \texttt{SÖYLE: T'yi (7,3)'e taşıdım. Şimdi sırada D var, D'yi (3,4)'e taşıyorum.} \\
            }
        }
    }
    & & \\ \\

    \theutterance \stepcounter{utterance}  
    & & \multicolumn{4}{p{0.6\linewidth}}{
        \cellcolor[rgb]{0.9,0.9,0.9}{
            \makecell[{{p{\linewidth}}}]{
                \texttt{\tiny{[P1$\langle$GM]}}
                \texttt{T’ başarıyla (7,3) koordinatlarına taşındı. Güncellenmiş ızgaran şu şekilde:} \\
\texttt{```} \\
\texttt{ 1234567} \\
\texttt{╔═╤══╤══╗} \\
\texttt{║◌│◌◌│◌◌║ 1} \\
\texttt{║◌└┬─┴──╢ 2} \\
\texttt{║◌◌│O◌◌T║ 3} \\
\texttt{╟─┬┘◌P┌─╢ 4} \\
\texttt{║◌│AD◌│◌║ 5} \\
\texttt{╟─┼─┬─┤M║ 6} \\
\texttt{║◌│S│◌│◌║ 7} \\
\texttt{╚═╧═╧═╧═╝} \\
\texttt{```} \\
\texttt{Şu anda en fazla 28 turdan 12. turu oynuyorsunuz.} \\
\texttt{Toplamda 16 cezanın 9 kadarını biriktirdiniz.Diğer oyuncu şu mesajı gönderdi:} \\
\\ 
\\ 
\texttt{"T'yi (7,3)'e taşıdım. Şimdi sırada D var, D'yi (3,4)'e taşıyorum."} \\
\\ 
\\ 
\texttt{Bir sonraki komutun nedir?} \\
            }
        }
    }
    & \\ \\

    \theutterance \stepcounter{utterance}  
    & \multicolumn{4}{p{0.6\linewidth}}{
        \cellcolor[rgb]{0.9,0.9,0.9}{
            \makecell[{{p{\linewidth}}}]{
                \texttt{\tiny{[P1$\rangle$GM]}}
                \texttt{HAMLE: D, (3,4)} \\
            }
        }
    }
    & & \\ \\

    \theutterance \stepcounter{utterance}  
    & & & \multicolumn{2}{p{0.3\linewidth}}{
        \cellcolor[rgb]{0.9,0.9,0.9}{
            \makecell[{{p{\linewidth}}}]{
                \texttt{\tiny{[GM$|$GM]}}
                \texttt{HAMLE: D, (3,4)} \\
            }
        }
    }
    & & \\ \\

    \theutterance \stepcounter{utterance}  
    & & & \multicolumn{2}{p{0.3\linewidth}}{
        \cellcolor[rgb]{0.9,0.9,0.9}{
            \makecell[{{p{\linewidth}}}]{
                \texttt{\tiny{[GM$|$GM]}}
                \texttt{Ceza: (3,4) boş değil; "┘" içeriyor.} \\
\\ 
\texttt{Toplamda 16 cezanın 10 kadarını biriktirdiniz.Lütfen tekrar deneyin!} \\
            }
        }
    }
    & & \\ \\

    \theutterance \stepcounter{utterance}  
    & & \multicolumn{4}{p{0.6\linewidth}}{
        \cellcolor[rgb]{0.9,0.9,0.9}{
            \makecell[{{p{\linewidth}}}]{
                \texttt{\tiny{[P1$\langle$GM]}}
                \texttt{Ceza: (3,4) boş değil; "┘" içeriyor.} \\
\\ 
\texttt{Toplamda 16 cezanın 10 kadarını biriktirdiniz.Lütfen tekrar deneyin!} \\
            }
        }
    }
    & \\ \\

    \theutterance \stepcounter{utterance}  
    & \multicolumn{4}{p{0.6\linewidth}}{
        \cellcolor[rgb]{0.9,0.9,0.9}{
            \makecell[{{p{\linewidth}}}]{
                \texttt{\tiny{[P1$\rangle$GM]}}
                \texttt{SÖYLE: (3,4) benim tahtamda ┘ sembolü var, boş değil. D için alternatif olarak (1,4) veya (2,4) boş. Hangisini seçelim? A hedefi (2,4) olduğu için D'yi (1,4) yapabiliriz. Sende (1,4) boş mu?} \\
            }
        }
    }
    & & \\ \\

    \theutterance \stepcounter{utterance}  
    & & & \multicolumn{2}{p{0.3\linewidth}}{
        \cellcolor[rgb]{0.9,0.9,0.9}{
            \makecell[{{p{\linewidth}}}]{
                \texttt{\tiny{[GM$|$GM]}}
                \texttt{SÖYLE: (3,4) benim tahtamda ┘ sembolü var, boş değil. D için alternatif olarak (1,4) veya (2,4) boş. Hangisini seçelim? A hedefi (2,4) olduğu için D'yi (1,4) yapabiliriz. Sende (1,4) boş mu?} \\
            }
        }
    }
    & & \\ \\

    \theutterance \stepcounter{utterance}  
    & & \multicolumn{4}{p{0.6\linewidth}}{
        \cellcolor[rgb]{0.9,0.9,0.9}{
            \makecell[{{p{\linewidth}}}]{
                \texttt{\tiny{[GM$\rangle$P2]}}
                \texttt{Mesajın diğer oyuncuya iletildi.} \\
\\ 
\\ 
\texttt{Şu anda en fazla 28 turdan 13. turu oynuyorsunuz.} \\
\texttt{Toplamda 16 cezanın 10 kadarını biriktirdiniz.Diğer oyuncu şu mesajı gönderdi:} \\
\\ 
\\ 
\texttt{"(3,4) benim tahtamda ┘ sembolü var, boş değil. D için alternatif olarak (1,4) veya (2,4) boş. Hangisini seçelim? A hedefi (2,4) olduğu için D'yi (1,4) yapabiliriz. Sende (1,4) boş mu?"} \\
\\ 
\\ 
\texttt{Bir sonraki komutun nedir?} \\
            }
        }
    }
    & \\ \\

    \theutterance \stepcounter{utterance}  
    & & & \multicolumn{4}{p{0.6\linewidth}}{
        \cellcolor[rgb]{0.9,0.9,0.9}{
            \makecell[{{p{\linewidth}}}]{
                \texttt{\tiny{[GM$\langle$P2]}}
                \texttt{SÖYLE: (1,4) benim tahtamda boş. D için (1,4) uygun. Güncel hedefler: P(5,4), T(7,3), D(1,4), A(2,4), S(4,6), M(5,6). D'yi (1,4)'e taşıyorum.} \\
            }
        }
    }
     \\ \\

    \theutterance \stepcounter{utterance}  
    & & & \multicolumn{2}{p{0.3\linewidth}}{
        \cellcolor[rgb]{0.9,0.9,0.9}{
            \makecell[{{p{\linewidth}}}]{
                \texttt{\tiny{[GM$|$GM]}}
                \texttt{SÖYLE: (1,4) benim tahtamda boş. D için (1,4) uygun. Güncel hedefler: P(5,4), T(7,3), D(1,4), A(2,4), S(4,6), M(5,6). D'yi (1,4)'e taşıyorum.} \\
            }
        }
    }
    & & \\ \\

    \theutterance \stepcounter{utterance}  
    & & \multicolumn{4}{p{0.6\linewidth}}{
        \cellcolor[rgb]{0.9,0.9,0.9}{
            \makecell[{{p{\linewidth}}}]{
                \texttt{\tiny{[P1$\langle$GM]}}
                \texttt{Mesajın diğer oyuncuya iletildi.} \\
\\ 
\\ 
\texttt{Şu anda en fazla 28 turdan 13. turu oynuyorsunuz.} \\
\texttt{Toplamda 16 cezanın 10 kadarını biriktirdiniz.Diğer oyuncu şu mesajı gönderdi:} \\
\\ 
\\ 
\texttt{"(1,4) benim tahtamda boş. D için (1,4) uygun. Güncel hedefler: P(5,4), T(7,3), D(1,4), A(2,4), S(4,6), M(5,6). D'yi (1,4)'e taşıyorum."} \\
\\ 
\\ 
\texttt{Bir sonraki komutun nedir?} \\
            }
        }
    }
    & \\ \\

    \theutterance \stepcounter{utterance}  
    & \multicolumn{4}{p{0.6\linewidth}}{
        \cellcolor[rgb]{0.9,0.9,0.9}{
            \makecell[{{p{\linewidth}}}]{
                \texttt{\tiny{[P1$\rangle$GM]}}
                \texttt{HAMLE: D, (1,4)} \\
            }
        }
    }
    & & \\ \\

    \theutterance \stepcounter{utterance}  
    & & & \multicolumn{2}{p{0.3\linewidth}}{
        \cellcolor[rgb]{0.9,0.9,0.9}{
            \makecell[{{p{\linewidth}}}]{
                \texttt{\tiny{[GM$|$GM]}}
                \texttt{HAMLE: D, (1,4)} \\
            }
        }
    }
    & & \\ \\

    \theutterance \stepcounter{utterance}  
    & & & \multicolumn{2}{p{0.3\linewidth}}{
        \cellcolor[rgb]{0.9,0.9,0.9}{
            \makecell[{{p{\linewidth}}}]{
                \texttt{\tiny{[GM$|$GM]}}
                \texttt{Ceza: (1,4) boş değil; "─" içeriyor.} \\
\\ 
\texttt{Toplamda 16 cezanın 11 kadarını biriktirdiniz.Lütfen tekrar deneyin!} \\
            }
        }
    }
    & & \\ \\

    \theutterance \stepcounter{utterance}  
    & & \multicolumn{4}{p{0.6\linewidth}}{
        \cellcolor[rgb]{0.9,0.9,0.9}{
            \makecell[{{p{\linewidth}}}]{
                \texttt{\tiny{[P1$\langle$GM]}}
                \texttt{Ceza: (1,4) boş değil; "─" içeriyor.} \\
\\ 
\texttt{Toplamda 16 cezanın 11 kadarını biriktirdiniz.Lütfen tekrar deneyin!} \\
            }
        }
    }
    & \\ \\

    \theutterance \stepcounter{utterance}  
    & \multicolumn{4}{p{0.6\linewidth}}{
        \cellcolor[rgb]{0.9,0.9,0.9}{
            \makecell[{{p{\linewidth}}}]{
                \texttt{\tiny{[P1$\rangle$GM]}}
                \texttt{SÖYLE: (1,4) de boş değil, ─ var. D için (2,3) veya (5,3) boş. Sende hangisi boş? A'yı da başka bir yere taşımamız gerekebilir çünkü (2,4) de └ var mı kontrol edeyim.} \\
            }
        }
    }
    & & \\ \\

    \theutterance \stepcounter{utterance}  
    & & & \multicolumn{2}{p{0.3\linewidth}}{
        \cellcolor[rgb]{0.9,0.9,0.9}{
            \makecell[{{p{\linewidth}}}]{
                \texttt{\tiny{[GM$|$GM]}}
                \texttt{SÖYLE: (1,4) de boş değil, ─ var. D için (2,3) veya (5,3) boş. Sende hangisi boş? A'yı da başka bir yere taşımamız gerekebilir çünkü (2,4) de └ var mı kontrol edeyim.} \\
            }
        }
    }
    & & \\ \\

    \theutterance \stepcounter{utterance}  
    & & \multicolumn{4}{p{0.6\linewidth}}{
        \cellcolor[rgb]{0.9,0.9,0.9}{
            \makecell[{{p{\linewidth}}}]{
                \texttt{\tiny{[GM$\rangle$P2]}}
                \texttt{Mesajın diğer oyuncuya iletildi.} \\
\\ 
\\ 
\texttt{Şu anda en fazla 28 turdan 14. turu oynuyorsunuz.} \\
\texttt{Toplamda 16 cezanın 11 kadarını biriktirdiniz.Diğer oyuncu şu mesajı gönderdi:} \\
\\ 
\\ 
\texttt{"(1,4) de boş değil, ─ var. D için (2,3) veya (5,3) boş. Sende hangisi boş? A'yı da başka bir yere taşımamız gerekebilir çünkü (2,4) de └ var mı kontrol edeyim."} \\
\\ 
\\ 
\texttt{Bir sonraki komutun nedir?} \\
            }
        }
    }
    & \\ \\

    \theutterance \stepcounter{utterance}  
    & & & \multicolumn{4}{p{0.6\linewidth}}{
        \cellcolor[rgb]{0.9,0.9,0.9}{
            \makecell[{{p{\linewidth}}}]{
                \texttt{\tiny{[GM$\langle$P2]}}
                \texttt{SÖYLE: (2,3) benim tahtamda boş, (5,3) ise D'nin şu anki konumu. D için (2,3) hedefini kullanalım. A için de (2,4) boş mu kontrol et, değilse yeni hedef belirleyelim. Bu arada D'yi (2,3)'e taşıyorum.} \\
            }
        }
    }
     \\ \\

    \theutterance \stepcounter{utterance}  
    & & & \multicolumn{2}{p{0.3\linewidth}}{
        \cellcolor[rgb]{0.9,0.9,0.9}{
            \makecell[{{p{\linewidth}}}]{
                \texttt{\tiny{[GM$|$GM]}}
                \texttt{SÖYLE: (2,3) benim tahtamda boş, (5,3) ise D'nin şu anki konumu. D için (2,3) hedefini kullanalım. A için de (2,4) boş mu kontrol et, değilse yeni hedef belirleyelim. Bu arada D'yi (2,3)'e taşıyorum.} \\
            }
        }
    }
    & & \\ \\

    \theutterance \stepcounter{utterance}  
    & & \multicolumn{4}{p{0.6\linewidth}}{
        \cellcolor[rgb]{0.9,0.9,0.9}{
            \makecell[{{p{\linewidth}}}]{
                \texttt{\tiny{[P1$\langle$GM]}}
                \texttt{Mesajın diğer oyuncuya iletildi.} \\
\\ 
\\ 
\texttt{Şu anda en fazla 28 turdan 14. turu oynuyorsunuz.} \\
\texttt{Toplamda 16 cezanın 11 kadarını biriktirdiniz.Diğer oyuncu şu mesajı gönderdi:} \\
\\ 
\\ 
\texttt{"(2,3) benim tahtamda boş, (5,3) ise D'nin şu anki konumu. D için (2,3) hedefini kullanalım. A için de (2,4) boş mu kontrol et, değilse yeni hedef belirleyelim. Bu arada D'yi (2,3)'e taşıyorum."} \\
\\ 
\\ 
\texttt{Bir sonraki komutun nedir?} \\
            }
        }
    }
    & \\ \\

    \theutterance \stepcounter{utterance}  
    & \multicolumn{4}{p{0.6\linewidth}}{
        \cellcolor[rgb]{0.9,0.9,0.9}{
            \makecell[{{p{\linewidth}}}]{
                \texttt{\tiny{[P1$\rangle$GM]}}
                \texttt{SÖYLE: (2,3) benim tahtamda boş, D için uygun. D'yi (2,3)'e taşıyorum. A için (2,4) boş mu diye kontrol edeyim, değilse (1,3) veya (6,3) önereceğim.} \\
            }
        }
    }
    & & \\ \\

    \theutterance \stepcounter{utterance}  
    & & & \multicolumn{2}{p{0.3\linewidth}}{
        \cellcolor[rgb]{0.9,0.9,0.9}{
            \makecell[{{p{\linewidth}}}]{
                \texttt{\tiny{[GM$|$GM]}}
                \texttt{SÖYLE: (2,3) benim tahtamda boş, D için uygun. D'yi (2,3)'e taşıyorum. A için (2,4) boş mu diye kontrol edeyim, değilse (1,3) veya (6,3) önereceğim.} \\
            }
        }
    }
    & & \\ \\

    \theutterance \stepcounter{utterance}  
    & & \multicolumn{4}{p{0.6\linewidth}}{
        \cellcolor[rgb]{0.9,0.9,0.9}{
            \makecell[{{p{\linewidth}}}]{
                \texttt{\tiny{[GM$\rangle$P2]}}
                \texttt{Mesajın diğer oyuncuya iletildi.} \\
\\ 
\\ 
\texttt{Şu anda en fazla 28 turdan 15. turu oynuyorsunuz.} \\
\texttt{Toplamda 16 cezanın 11 kadarını biriktirdiniz.Diğer oyuncu şu mesajı gönderdi:} \\
\\ 
\\ 
\texttt{"(2,3) benim tahtamda boş, D için uygun. D'yi (2,3)'e taşıyorum. A için (2,4) boş mu diye kontrol edeyim, değilse (1,3) veya (6,3) önereceğim."} \\
\\ 
\\ 
\texttt{Bir sonraki komutun nedir?} \\
            }
        }
    }
    & \\ \\

    \theutterance \stepcounter{utterance}  
    & & & \multicolumn{4}{p{0.6\linewidth}}{
        \cellcolor[rgb]{0.9,0.9,0.9}{
            \makecell[{{p{\linewidth}}}]{
                \texttt{\tiny{[GM$\langle$P2]}}
                \texttt{HAMLE: D, (2,3)} \\
            }
        }
    }
     \\ \\

    \theutterance \stepcounter{utterance}  
    & & & \multicolumn{2}{p{0.3\linewidth}}{
        \cellcolor[rgb]{0.9,0.9,0.9}{
            \makecell[{{p{\linewidth}}}]{
                \texttt{\tiny{[GM$|$GM]}}
                \texttt{HAMLE: D, (2,3)} \\
            }
        }
    }
    & & \\ \\

    \theutterance \stepcounter{utterance}  
    & & & \multicolumn{2}{p{0.3\linewidth}}{
        \cellcolor[rgb]{0.9,0.9,0.9}{
            \makecell[{{p{\linewidth}}}]{
                \texttt{\tiny{[GM$|$GM]}}
                \texttt{D moved to (2, 3)} \\
            }
        }
    }
    & & \\ \\

    \theutterance \stepcounter{utterance}  
    & & \multicolumn{4}{p{0.6\linewidth}}{
        \cellcolor[rgb]{0.9,0.9,0.9}{
            \makecell[{{p{\linewidth}}}]{
                \texttt{\tiny{[P1$\langle$GM]}}
                \texttt{Mesajın diğer oyuncuya iletildi.} \\
\\ 
\\ 
\texttt{Şu anda en fazla 28 turdan 15. turu oynuyorsunuz.} \\
\texttt{Toplamda 16 cezanın 11 kadarını biriktirdiniz.Diğer oyuncu kendi ızgarasında bir nesneyi hareket ettirdi.} \\
\\ 
\\ 
\texttt{Bir sonraki komutun nedir?} \\
            }
        }
    }
    & \\ \\

    \theutterance \stepcounter{utterance}  
    & \multicolumn{4}{p{0.6\linewidth}}{
        \cellcolor[rgb]{0.9,0.9,0.9}{
            \makecell[{{p{\linewidth}}}]{
                \texttt{\tiny{[P1$\rangle$GM]}}
                \texttt{HAMLE: D, (2,3)} \\
            }
        }
    }
    & & \\ \\

    \theutterance \stepcounter{utterance}  
    & & & \multicolumn{2}{p{0.3\linewidth}}{
        \cellcolor[rgb]{0.9,0.9,0.9}{
            \makecell[{{p{\linewidth}}}]{
                \texttt{\tiny{[GM$|$GM]}}
                \texttt{HAMLE: D, (2,3)} \\
            }
        }
    }
    & & \\ \\

    \theutterance \stepcounter{utterance}  
    & & & \multicolumn{2}{p{0.3\linewidth}}{
        \cellcolor[rgb]{0.9,0.9,0.9}{
            \makecell[{{p{\linewidth}}}]{
                \texttt{\tiny{[GM$|$GM]}}
                \texttt{D moved to (2, 3)} \\
            }
        }
    }
    & & \\ \\

    \theutterance \stepcounter{utterance}  
    & & \multicolumn{4}{p{0.6\linewidth}}{
        \cellcolor[rgb]{0.9,0.9,0.9}{
            \makecell[{{p{\linewidth}}}]{
                \texttt{\tiny{[GM$\rangle$P2]}}
                \texttt{D’ başarıyla (2,3) koordinatlarına taşındı. Güncellenmiş ızgaran şu şekilde:} \\
\texttt{```} \\
\texttt{ 1234567} \\
\texttt{╔═╤══╤══╗} \\
\texttt{║◌│O◌│◌◌║ 1} \\
\texttt{║◌└┬─┴──╢ 2} \\
\texttt{║◌D│◌◌◌T║ 3} \\
\texttt{╟─┬┘◌P┌─╢ 4} \\
\texttt{║◌│◌◌◌│A║ 5} \\
\texttt{╟─┼─┬─┤S║ 6} \\
\texttt{║M│◌│◌│◌║ 7} \\
\texttt{╚═╧═╧═╧═╝} \\
\texttt{```} \\
\texttt{Şu anda en fazla 28 turdan 16. turu oynuyorsunuz.} \\
\texttt{Toplamda 16 cezanın 11 kadarını biriktirdiniz.Diğer oyuncu kendi ızgarasında bir nesneyi hareket ettirdi.} \\
\\ 
\\ 
\texttt{Bir sonraki komutun nedir?} \\
            }
        }
    }
    & \\ \\

    \theutterance \stepcounter{utterance}  
    & & & \multicolumn{4}{p{0.6\linewidth}}{
        \cellcolor[rgb]{0.9,0.9,0.9}{
            \makecell[{{p{\linewidth}}}]{
                \texttt{\tiny{[GM$\langle$P2]}}
                \texttt{SÖYLE: D'yi (2,3)'e taşıdım. A için (2,4) boş mu diye kontrol ettin mi? Değilse alternatif öner, ben de kontrol edeyim.} \\
            }
        }
    }
     \\ \\

    \theutterance \stepcounter{utterance}  
    & & & \multicolumn{2}{p{0.3\linewidth}}{
        \cellcolor[rgb]{0.9,0.9,0.9}{
            \makecell[{{p{\linewidth}}}]{
                \texttt{\tiny{[GM$|$GM]}}
                \texttt{SÖYLE: D'yi (2,3)'e taşıdım. A için (2,4) boş mu diye kontrol ettin mi? Değilse alternatif öner, ben de kontrol edeyim.} \\
            }
        }
    }
    & & \\ \\

    \theutterance \stepcounter{utterance}  
    & & \multicolumn{4}{p{0.6\linewidth}}{
        \cellcolor[rgb]{0.9,0.9,0.9}{
            \makecell[{{p{\linewidth}}}]{
                \texttt{\tiny{[P1$\langle$GM]}}
                \texttt{D’ başarıyla (2,3) koordinatlarına taşındı. Güncellenmiş ızgaran şu şekilde:} \\
\texttt{```} \\
\texttt{ 1234567} \\
\texttt{╔═╤══╤══╗} \\
\texttt{║◌│◌◌│◌◌║ 1} \\
\texttt{║◌└┬─┴──╢ 2} \\
\texttt{║◌D│O◌◌T║ 3} \\
\texttt{╟─┬┘◌P┌─╢ 4} \\
\texttt{║◌│A◌◌│◌║ 5} \\
\texttt{╟─┼─┬─┤M║ 6} \\
\texttt{║◌│S│◌│◌║ 7} \\
\texttt{╚═╧═╧═╧═╝} \\
\texttt{```} \\
\texttt{Şu anda en fazla 28 turdan 16. turu oynuyorsunuz.} \\
\texttt{Toplamda 16 cezanın 11 kadarını biriktirdiniz.Diğer oyuncu şu mesajı gönderdi:} \\
\\ 
\\ 
\texttt{"D'yi (2,3)'e taşıdım. A için (2,4) boş mu diye kontrol ettin mi? Değilse alternatif öner, ben de kontrol edeyim."} \\
\\ 
\\ 
\texttt{Bir sonraki komutun nedir?} \\
            }
        }
    }
    & \\ \\

    \theutterance \stepcounter{utterance}  
    & \multicolumn{4}{p{0.6\linewidth}}{
        \cellcolor[rgb]{0.9,0.9,0.9}{
            \makecell[{{p{\linewidth}}}]{
                \texttt{\tiny{[P1$\rangle$GM]}}
                \texttt{SÖYLE: (2,4) benim tahtamda └ var, boş değil. A için alternatif olarak (6,3) veya (1,3) boş. Hangisi sende boş? (6,3) T'nin yanında olur, (1,3) ise D'nin yanında olur. Hangisini tercih edersin?} \\
            }
        }
    }
    & & \\ \\

    \theutterance \stepcounter{utterance}  
    & & & \multicolumn{2}{p{0.3\linewidth}}{
        \cellcolor[rgb]{0.9,0.9,0.9}{
            \makecell[{{p{\linewidth}}}]{
                \texttt{\tiny{[GM$|$GM]}}
                \texttt{SÖYLE: (2,4) benim tahtamda └ var, boş değil. A için alternatif olarak (6,3) veya (1,3) boş. Hangisi sende boş? (6,3) T'nin yanında olur, (1,3) ise D'nin yanında olur. Hangisini tercih edersin?} \\
            }
        }
    }
    & & \\ \\

    \theutterance \stepcounter{utterance}  
    & & \multicolumn{4}{p{0.6\linewidth}}{
        \cellcolor[rgb]{0.9,0.9,0.9}{
            \makecell[{{p{\linewidth}}}]{
                \texttt{\tiny{[GM$\rangle$P2]}}
                \texttt{Mesajın diğer oyuncuya iletildi.} \\
\\ 
\\ 
\texttt{Şu anda en fazla 28 turdan 17. turu oynuyorsunuz.} \\
\texttt{Toplamda 16 cezanın 11 kadarını biriktirdiniz.Diğer oyuncu şu mesajı gönderdi:} \\
\\ 
\\ 
\texttt{"(2,4) benim tahtamda └ var, boş değil. A için alternatif olarak (6,3) veya (1,3) boş. Hangisi sende boş? (6,3) T'nin yanında olur, (1,3) ise D'nin yanında olur. Hangisini tercih edersin?"} \\
\\ 
\\ 
\texttt{Bir sonraki komutun nedir?} \\
            }
        }
    }
    & \\ \\

    \theutterance \stepcounter{utterance}  
    & & & \multicolumn{4}{p{0.6\linewidth}}{
        \cellcolor[rgb]{0.9,0.9,0.9}{
            \makecell[{{p{\linewidth}}}]{
                \texttt{\tiny{[GM$\langle$P2]}}
                \texttt{SÖYLE: (1,3) benim tahtamda boş, (6,3) de boş. A için (1,3) uygun, D'nin yanında olsun. Güncel hedefler: P(5,4), T(7,3), D(2,3), A(1,3), S(4,6), M(5,6). A'yı (1,3)'e taşıyorum.} \\
            }
        }
    }
     \\ \\

    \theutterance \stepcounter{utterance}  
    & & & \multicolumn{2}{p{0.3\linewidth}}{
        \cellcolor[rgb]{0.9,0.9,0.9}{
            \makecell[{{p{\linewidth}}}]{
                \texttt{\tiny{[GM$|$GM]}}
                \texttt{SÖYLE: (1,3) benim tahtamda boş, (6,3) de boş. A için (1,3) uygun, D'nin yanında olsun. Güncel hedefler: P(5,4), T(7,3), D(2,3), A(1,3), S(4,6), M(5,6). A'yı (1,3)'e taşıyorum.} \\
            }
        }
    }
    & & \\ \\

    \theutterance \stepcounter{utterance}  
    & & \multicolumn{4}{p{0.6\linewidth}}{
        \cellcolor[rgb]{0.9,0.9,0.9}{
            \makecell[{{p{\linewidth}}}]{
                \texttt{\tiny{[P1$\langle$GM]}}
                \texttt{Mesajın diğer oyuncuya iletildi.} \\
\\ 
\\ 
\texttt{Şu anda en fazla 28 turdan 17. turu oynuyorsunuz.} \\
\texttt{Toplamda 16 cezanın 11 kadarını biriktirdiniz.Diğer oyuncu şu mesajı gönderdi:} \\
\\ 
\\ 
\texttt{"(1,3) benim tahtamda boş, (6,3) de boş. A için (1,3) uygun, D'nin yanında olsun. Güncel hedefler: P(5,4), T(7,3), D(2,3), A(1,3), S(4,6), M(5,6). A'yı (1,3)'e taşıyorum."} \\
\\ 
\\ 
\texttt{Bir sonraki komutun nedir?} \\
            }
        }
    }
    & \\ \\

    \theutterance \stepcounter{utterance}  
    & \multicolumn{4}{p{0.6\linewidth}}{
        \cellcolor[rgb]{0.9,0.9,0.9}{
            \makecell[{{p{\linewidth}}}]{
                \texttt{\tiny{[P1$\rangle$GM]}}
                \texttt{HAMLE: A, (1,3)} \\
            }
        }
    }
    & & \\ \\

    \theutterance \stepcounter{utterance}  
    & & & \multicolumn{2}{p{0.3\linewidth}}{
        \cellcolor[rgb]{0.9,0.9,0.9}{
            \makecell[{{p{\linewidth}}}]{
                \texttt{\tiny{[GM$|$GM]}}
                \texttt{HAMLE: A, (1,3)} \\
            }
        }
    }
    & & \\ \\

    \theutterance \stepcounter{utterance}  
    & & & \multicolumn{2}{p{0.3\linewidth}}{
        \cellcolor[rgb]{0.9,0.9,0.9}{
            \makecell[{{p{\linewidth}}}]{
                \texttt{\tiny{[GM$|$GM]}}
                \texttt{A moved to (1, 3)} \\
            }
        }
    }
    & & \\ \\

    \theutterance \stepcounter{utterance}  
    & & \multicolumn{4}{p{0.6\linewidth}}{
        \cellcolor[rgb]{0.9,0.9,0.9}{
            \makecell[{{p{\linewidth}}}]{
                \texttt{\tiny{[GM$\rangle$P2]}}
                \texttt{Mesajın diğer oyuncuya iletildi.} \\
\\ 
\\ 
\texttt{Şu anda en fazla 28 turdan 18. turu oynuyorsunuz.} \\
\texttt{Toplamda 16 cezanın 11 kadarını biriktirdiniz.Diğer oyuncu kendi ızgarasında bir nesneyi hareket ettirdi.} \\
\\ 
\\ 
\texttt{Bir sonraki komutun nedir?} \\
            }
        }
    }
    & \\ \\

    \theutterance \stepcounter{utterance}  
    & & & \multicolumn{4}{p{0.6\linewidth}}{
        \cellcolor[rgb]{0.9,0.9,0.9}{
            \makecell[{{p{\linewidth}}}]{
                \texttt{\tiny{[GM$\langle$P2]}}
                \texttt{HAMLE: A, (1,3)} \\
            }
        }
    }
     \\ \\

    \theutterance \stepcounter{utterance}  
    & & & \multicolumn{2}{p{0.3\linewidth}}{
        \cellcolor[rgb]{0.9,0.9,0.9}{
            \makecell[{{p{\linewidth}}}]{
                \texttt{\tiny{[GM$|$GM]}}
                \texttt{HAMLE: A, (1,3)} \\
            }
        }
    }
    & & \\ \\

    \theutterance \stepcounter{utterance}  
    & & & \multicolumn{2}{p{0.3\linewidth}}{
        \cellcolor[rgb]{0.9,0.9,0.9}{
            \makecell[{{p{\linewidth}}}]{
                \texttt{\tiny{[GM$|$GM]}}
                \texttt{A moved to (1, 3)} \\
            }
        }
    }
    & & \\ \\

    \theutterance \stepcounter{utterance}  
    & & \multicolumn{4}{p{0.6\linewidth}}{
        \cellcolor[rgb]{0.9,0.9,0.9}{
            \makecell[{{p{\linewidth}}}]{
                \texttt{\tiny{[P1$\langle$GM]}}
                \texttt{A’ başarıyla (1,3) koordinatlarına taşındı. Güncellenmiş ızgaran şu şekilde:} \\
\texttt{```} \\
\texttt{ 1234567} \\
\texttt{╔═╤══╤══╗} \\
\texttt{║◌│◌◌│◌◌║ 1} \\
\texttt{║◌└┬─┴──╢ 2} \\
\texttt{║AD│O◌◌T║ 3} \\
\texttt{╟─┬┘◌P┌─╢ 4} \\
\texttt{║◌│◌◌◌│◌║ 5} \\
\texttt{╟─┼─┬─┤M║ 6} \\
\texttt{║◌│S│◌│◌║ 7} \\
\texttt{╚═╧═╧═╧═╝} \\
\texttt{```} \\
\texttt{Şu anda en fazla 28 turdan 18. turu oynuyorsunuz.} \\
\texttt{Toplamda 16 cezanın 11 kadarını biriktirdiniz.Diğer oyuncu kendi ızgarasında bir nesneyi hareket ettirdi.} \\
\\ 
\\ 
\texttt{Bir sonraki komutun nedir?} \\
            }
        }
    }
    & \\ \\

    \theutterance \stepcounter{utterance}  
    & \multicolumn{4}{p{0.6\linewidth}}{
        \cellcolor[rgb]{0.9,0.9,0.9}{
            \makecell[{{p{\linewidth}}}]{
                \texttt{\tiny{[P1$\rangle$GM]}}
                \texttt{SÖYLE: A'yı (1,3)'e taşıdım. Şimdi sırada S var. S için hedef (4,6) hâlâ uygun mu, sende boş mu? Boşsa S'yi (4,6)'ya taşıyalım.} \\
            }
        }
    }
    & & \\ \\

    \theutterance \stepcounter{utterance}  
    & & & \multicolumn{2}{p{0.3\linewidth}}{
        \cellcolor[rgb]{0.9,0.9,0.9}{
            \makecell[{{p{\linewidth}}}]{
                \texttt{\tiny{[GM$|$GM]}}
                \texttt{SÖYLE: A'yı (1,3)'e taşıdım. Şimdi sırada S var. S için hedef (4,6) hâlâ uygun mu, sende boş mu? Boşsa S'yi (4,6)'ya taşıyalım.} \\
            }
        }
    }
    & & \\ \\

    \theutterance \stepcounter{utterance}  
    & & \multicolumn{4}{p{0.6\linewidth}}{
        \cellcolor[rgb]{0.9,0.9,0.9}{
            \makecell[{{p{\linewidth}}}]{
                \texttt{\tiny{[GM$\rangle$P2]}}
                \texttt{A’ başarıyla (1,3) koordinatlarına taşındı. Güncellenmiş ızgaran şu şekilde:} \\
\texttt{```} \\
\texttt{ 1234567} \\
\texttt{╔═╤══╤══╗} \\
\texttt{║◌│O◌│◌◌║ 1} \\
\texttt{║◌└┬─┴──╢ 2} \\
\texttt{║AD│◌◌◌T║ 3} \\
\texttt{╟─┬┘◌P┌─╢ 4} \\
\texttt{║◌│◌◌◌│◌║ 5} \\
\texttt{╟─┼─┬─┤S║ 6} \\
\texttt{║M│◌│◌│◌║ 7} \\
\texttt{╚═╧═╧═╧═╝} \\
\texttt{```} \\
\texttt{Şu anda en fazla 28 turdan 19. turu oynuyorsunuz.} \\
\texttt{Toplamda 16 cezanın 11 kadarını biriktirdiniz.Diğer oyuncu şu mesajı gönderdi:} \\
\\ 
\\ 
\texttt{"A'yı (1,3)'e taşıdım. Şimdi sırada S var. S için hedef (4,6) hâlâ uygun mu, sende boş mu? Boşsa S'yi (4,6)'ya taşıyalım."} \\
\\ 
\\ 
\texttt{Bir sonraki komutun nedir?} \\
            }
        }
    }
    & \\ \\

    \theutterance \stepcounter{utterance}  
    & & & \multicolumn{4}{p{0.6\linewidth}}{
        \cellcolor[rgb]{0.9,0.9,0.9}{
            \makecell[{{p{\linewidth}}}]{
                \texttt{\tiny{[GM$\langle$P2]}}
                \texttt{SÖYLE: (4,6) benim tahtamda boş. S için uygun. S'yi (4,6)'ya taşıyorum.} \\
            }
        }
    }
     \\ \\

    \theutterance \stepcounter{utterance}  
    & & & \multicolumn{2}{p{0.3\linewidth}}{
        \cellcolor[rgb]{0.9,0.9,0.9}{
            \makecell[{{p{\linewidth}}}]{
                \texttt{\tiny{[GM$|$GM]}}
                \texttt{SÖYLE: (4,6) benim tahtamda boş. S için uygun. S'yi (4,6)'ya taşıyorum.} \\
            }
        }
    }
    & & \\ \\

    \theutterance \stepcounter{utterance}  
    & & \multicolumn{4}{p{0.6\linewidth}}{
        \cellcolor[rgb]{0.9,0.9,0.9}{
            \makecell[{{p{\linewidth}}}]{
                \texttt{\tiny{[P1$\langle$GM]}}
                \texttt{Mesajın diğer oyuncuya iletildi.} \\
\\ 
\\ 
\texttt{Şu anda en fazla 28 turdan 19. turu oynuyorsunuz.} \\
\texttt{Toplamda 16 cezanın 11 kadarını biriktirdiniz.Diğer oyuncu şu mesajı gönderdi:} \\
\\ 
\\ 
\texttt{"(4,6) benim tahtamda boş. S için uygun. S'yi (4,6)'ya taşıyorum."} \\
\\ 
\\ 
\texttt{Bir sonraki komutun nedir?} \\
            }
        }
    }
    & \\ \\

    \theutterance \stepcounter{utterance}  
    & \multicolumn{4}{p{0.6\linewidth}}{
        \cellcolor[rgb]{0.9,0.9,0.9}{
            \makecell[{{p{\linewidth}}}]{
                \texttt{\tiny{[P1$\rangle$GM]}}
                \texttt{HAMLE: S, (4,6)} \\
            }
        }
    }
    & & \\ \\

    \theutterance \stepcounter{utterance}  
    & & & \multicolumn{2}{p{0.3\linewidth}}{
        \cellcolor[rgb]{0.9,0.9,0.9}{
            \makecell[{{p{\linewidth}}}]{
                \texttt{\tiny{[GM$|$GM]}}
                \texttt{HAMLE: S, (4,6)} \\
            }
        }
    }
    & & \\ \\

    \theutterance \stepcounter{utterance}  
    & & & \multicolumn{2}{p{0.3\linewidth}}{
        \cellcolor[rgb]{0.9,0.9,0.9}{
            \makecell[{{p{\linewidth}}}]{
                \texttt{\tiny{[GM$|$GM]}}
                \texttt{Ceza: (4,6) boş değil; "┬" içeriyor.} \\
\\ 
\texttt{Toplamda 16 cezanın 12 kadarını biriktirdiniz.Lütfen tekrar deneyin!} \\
            }
        }
    }
    & & \\ \\

    \theutterance \stepcounter{utterance}  
    & & \multicolumn{4}{p{0.6\linewidth}}{
        \cellcolor[rgb]{0.9,0.9,0.9}{
            \makecell[{{p{\linewidth}}}]{
                \texttt{\tiny{[P1$\langle$GM]}}
                \texttt{Ceza: (4,6) boş değil; "┬" içeriyor.} \\
\\ 
\texttt{Toplamda 16 cezanın 12 kadarını biriktirdiniz.Lütfen tekrar deneyin!} \\
            }
        }
    }
    & \\ \\

    \theutterance \stepcounter{utterance}  
    & \multicolumn{4}{p{0.6\linewidth}}{
        \cellcolor[rgb]{0.9,0.9,0.9}{
            \makecell[{{p{\linewidth}}}]{
                \texttt{\tiny{[P1$\rangle$GM]}}
                \texttt{SÖYLE: (4,6) benim tahtamda ┬ var, boş değil. S için alternatif olarak (3,6) veya (5,6) boş. Ama (5,6) M'nin hedefiydi. S için (3,6) uygun mu? M için de (5,6) veya başka bir yer belirleyelim.} \\
            }
        }
    }
    & & \\ \\

    \theutterance \stepcounter{utterance}  
    & & & \multicolumn{2}{p{0.3\linewidth}}{
        \cellcolor[rgb]{0.9,0.9,0.9}{
            \makecell[{{p{\linewidth}}}]{
                \texttt{\tiny{[GM$|$GM]}}
                \texttt{SÖYLE: (4,6) benim tahtamda ┬ var, boş değil. S için alternatif olarak (3,6) veya (5,6) boş. Ama (5,6) M'nin hedefiydi. S için (3,6) uygun mu? M için de (5,6) veya başka bir yer belirleyelim.} \\
            }
        }
    }
    & & \\ \\

    \theutterance \stepcounter{utterance}  
    & & \multicolumn{4}{p{0.6\linewidth}}{
        \cellcolor[rgb]{0.9,0.9,0.9}{
            \makecell[{{p{\linewidth}}}]{
                \texttt{\tiny{[GM$\rangle$P2]}}
                \texttt{Mesajın diğer oyuncuya iletildi.} \\
\\ 
\\ 
\texttt{Şu anda en fazla 28 turdan 20. turu oynuyorsunuz.} \\
\texttt{Toplamda 16 cezanın 12 kadarını biriktirdiniz.Diğer oyuncu şu mesajı gönderdi:} \\
\\ 
\\ 
\texttt{"(4,6) benim tahtamda ┬ var, boş değil. S için alternatif olarak (3,6) veya (5,6) boş. Ama (5,6) M'nin hedefiydi. S için (3,6) uygun mu? M için de (5,6) veya başka bir yer belirleyelim."} \\
\\ 
\\ 
\texttt{Bir sonraki komutun nedir?} \\
            }
        }
    }
    & \\ \\

    \theutterance \stepcounter{utterance}  
    & & & \multicolumn{4}{p{0.6\linewidth}}{
        \cellcolor[rgb]{0.9,0.9,0.9}{
            \makecell[{{p{\linewidth}}}]{
                \texttt{\tiny{[GM$\langle$P2]}}
                \texttt{SÖYLE: (3,6) benim tahtamda boş. S için (3,6) uygun. M için (5,6) da boş, hâlâ uygun. Güncel hedefler: P(5,4), T(7,3), D(2,3), A(1,3), S(3,6), M(5,6). S'yi (3,6)'ya taşıyorum.} \\
            }
        }
    }
     \\ \\

    \theutterance \stepcounter{utterance}  
    & & & \multicolumn{2}{p{0.3\linewidth}}{
        \cellcolor[rgb]{0.9,0.9,0.9}{
            \makecell[{{p{\linewidth}}}]{
                \texttt{\tiny{[GM$|$GM]}}
                \texttt{SÖYLE: (3,6) benim tahtamda boş. S için (3,6) uygun. M için (5,6) da boş, hâlâ uygun. Güncel hedefler: P(5,4), T(7,3), D(2,3), A(1,3), S(3,6), M(5,6). S'yi (3,6)'ya taşıyorum.} \\
            }
        }
    }
    & & \\ \\

    \theutterance \stepcounter{utterance}  
    & & \multicolumn{4}{p{0.6\linewidth}}{
        \cellcolor[rgb]{0.9,0.9,0.9}{
            \makecell[{{p{\linewidth}}}]{
                \texttt{\tiny{[P1$\langle$GM]}}
                \texttt{Mesajın diğer oyuncuya iletildi.} \\
\\ 
\\ 
\texttt{Şu anda en fazla 28 turdan 20. turu oynuyorsunuz.} \\
\texttt{Toplamda 16 cezanın 12 kadarını biriktirdiniz.Diğer oyuncu şu mesajı gönderdi:} \\
\\ 
\\ 
\texttt{"(3,6) benim tahtamda boş. S için (3,6) uygun. M için (5,6) da boş, hâlâ uygun. Güncel hedefler: P(5,4), T(7,3), D(2,3), A(1,3), S(3,6), M(5,6). S'yi (3,6)'ya taşıyorum."} \\
\\ 
\\ 
\texttt{Bir sonraki komutun nedir?} \\
            }
        }
    }
    & \\ \\

    \theutterance \stepcounter{utterance}  
    & \multicolumn{4}{p{0.6\linewidth}}{
        \cellcolor[rgb]{0.9,0.9,0.9}{
            \makecell[{{p{\linewidth}}}]{
                \texttt{\tiny{[P1$\rangle$GM]}}
                \texttt{HAMLE: S, (3,6)} \\
            }
        }
    }
    & & \\ \\

    \theutterance \stepcounter{utterance}  
    & & & \multicolumn{2}{p{0.3\linewidth}}{
        \cellcolor[rgb]{0.9,0.9,0.9}{
            \makecell[{{p{\linewidth}}}]{
                \texttt{\tiny{[GM$|$GM]}}
                \texttt{HAMLE: S, (3,6)} \\
            }
        }
    }
    & & \\ \\

    \theutterance \stepcounter{utterance}  
    & & & \multicolumn{2}{p{0.3\linewidth}}{
        \cellcolor[rgb]{0.9,0.9,0.9}{
            \makecell[{{p{\linewidth}}}]{
                \texttt{\tiny{[GM$|$GM]}}
                \texttt{Ceza: (3,6) boş değil; "─" içeriyor.} \\
\\ 
\texttt{Toplamda 16 cezanın 13 kadarını biriktirdiniz.Lütfen tekrar deneyin!} \\
            }
        }
    }
    & & \\ \\

    \theutterance \stepcounter{utterance}  
    & & \multicolumn{4}{p{0.6\linewidth}}{
        \cellcolor[rgb]{0.9,0.9,0.9}{
            \makecell[{{p{\linewidth}}}]{
                \texttt{\tiny{[P1$\langle$GM]}}
                \texttt{Ceza: (3,6) boş değil; "─" içeriyor.} \\
\\ 
\texttt{Toplamda 16 cezanın 13 kadarını biriktirdiniz.Lütfen tekrar deneyin!} \\
            }
        }
    }
    & \\ \\

    \theutterance \stepcounter{utterance}  
    & \multicolumn{4}{p{0.6\linewidth}}{
        \cellcolor[rgb]{0.9,0.9,0.9}{
            \makecell[{{p{\linewidth}}}]{
                \texttt{\tiny{[P1$\rangle$GM]}}
                \texttt{SÖYLE: (3,6) da boş değil, ─ var. S için (1,6) veya (2,6) veya (5,6) boş. (5,6) M'nin hedefiydi ama S'yi oraya koysak M için başka yer buluruz. S için (1,6) uygun mu sende?} \\
            }
        }
    }
    & & \\ \\

    \theutterance \stepcounter{utterance}  
    & & & \multicolumn{2}{p{0.3\linewidth}}{
        \cellcolor[rgb]{0.9,0.9,0.9}{
            \makecell[{{p{\linewidth}}}]{
                \texttt{\tiny{[GM$|$GM]}}
                \texttt{SÖYLE: (3,6) da boş değil, ─ var. S için (1,6) veya (2,6) veya (5,6) boş. (5,6) M'nin hedefiydi ama S'yi oraya koysak M için başka yer buluruz. S için (1,6) uygun mu sende?} \\
            }
        }
    }
    & & \\ \\

    \theutterance \stepcounter{utterance}  
    & & \multicolumn{4}{p{0.6\linewidth}}{
        \cellcolor[rgb]{0.9,0.9,0.9}{
            \makecell[{{p{\linewidth}}}]{
                \texttt{\tiny{[GM$\rangle$P2]}}
                \texttt{Mesajın diğer oyuncuya iletildi.} \\
\\ 
\\ 
\texttt{Şu anda en fazla 28 turdan 21. turu oynuyorsunuz.} \\
\texttt{Toplamda 16 cezanın 13 kadarını biriktirdiniz.Diğer oyuncu şu mesajı gönderdi:} \\
\\ 
\\ 
\texttt{"(3,6) da boş değil, ─ var. S için (1,6) veya (2,6) veya (5,6) boş. (5,6) M'nin hedefiydi ama S'yi oraya koysak M için başka yer buluruz. S için (1,6) uygun mu sende?"} \\
\\ 
\\ 
\texttt{Bir sonraki komutun nedir?} \\
            }
        }
    }
    & \\ \\

    \theutterance \stepcounter{utterance}  
    & & & \multicolumn{4}{p{0.6\linewidth}}{
        \cellcolor[rgb]{0.9,0.9,0.9}{
            \makecell[{{p{\linewidth}}}]{
                \texttt{\tiny{[GM$\langle$P2]}}
                \texttt{SÖYLE: (1,6) benim tahtamda boş. S için (1,6) uygun. M için (5,6) hâlâ boş, oraya koyalım. Güncel hedefler: P(5,4), T(7,3), D(2,3), A(1,3), S(1,6), M(5,6). S'yi (1,6)'ya taşıyorum.} \\
            }
        }
    }
     \\ \\

    \theutterance \stepcounter{utterance}  
    & & & \multicolumn{2}{p{0.3\linewidth}}{
        \cellcolor[rgb]{0.9,0.9,0.9}{
            \makecell[{{p{\linewidth}}}]{
                \texttt{\tiny{[GM$|$GM]}}
                \texttt{SÖYLE: (1,6) benim tahtamda boş. S için (1,6) uygun. M için (5,6) hâlâ boş, oraya koyalım. Güncel hedefler: P(5,4), T(7,3), D(2,3), A(1,3), S(1,6), M(5,6). S'yi (1,6)'ya taşıyorum.} \\
            }
        }
    }
    & & \\ \\

    \theutterance \stepcounter{utterance}  
    & & \multicolumn{4}{p{0.6\linewidth}}{
        \cellcolor[rgb]{0.9,0.9,0.9}{
            \makecell[{{p{\linewidth}}}]{
                \texttt{\tiny{[P1$\langle$GM]}}
                \texttt{Mesajın diğer oyuncuya iletildi.} \\
\\ 
\\ 
\texttt{Şu anda en fazla 28 turdan 21. turu oynuyorsunuz.} \\
\texttt{Toplamda 16 cezanın 13 kadarını biriktirdiniz.Diğer oyuncu şu mesajı gönderdi:} \\
\\ 
\\ 
\texttt{"(1,6) benim tahtamda boş. S için (1,6) uygun. M için (5,6) hâlâ boş, oraya koyalım. Güncel hedefler: P(5,4), T(7,3), D(2,3), A(1,3), S(1,6), M(5,6). S'yi (1,6)'ya taşıyorum."} \\
\\ 
\\ 
\texttt{Bir sonraki komutun nedir?} \\
            }
        }
    }
    & \\ \\

    \theutterance \stepcounter{utterance}  
    & \multicolumn{4}{p{0.6\linewidth}}{
        \cellcolor[rgb]{0.9,0.9,0.9}{
            \makecell[{{p{\linewidth}}}]{
                \texttt{\tiny{[P1$\rangle$GM]}}
                \texttt{HAMLE: S, (1,6)} \\
            }
        }
    }
    & & \\ \\

    \theutterance \stepcounter{utterance}  
    & & & \multicolumn{2}{p{0.3\linewidth}}{
        \cellcolor[rgb]{0.9,0.9,0.9}{
            \makecell[{{p{\linewidth}}}]{
                \texttt{\tiny{[GM$|$GM]}}
                \texttt{HAMLE: S, (1,6)} \\
            }
        }
    }
    & & \\ \\

    \theutterance \stepcounter{utterance}  
    & & & \multicolumn{2}{p{0.3\linewidth}}{
        \cellcolor[rgb]{0.9,0.9,0.9}{
            \makecell[{{p{\linewidth}}}]{
                \texttt{\tiny{[GM$|$GM]}}
                \texttt{Ceza: (1,6) boş değil; "─" içeriyor.} \\
\\ 
\texttt{Toplamda 16 cezanın 14 kadarını biriktirdiniz.Lütfen tekrar deneyin!} \\
            }
        }
    }
    & & \\ \\

    \theutterance \stepcounter{utterance}  
    & & \multicolumn{4}{p{0.6\linewidth}}{
        \cellcolor[rgb]{0.9,0.9,0.9}{
            \makecell[{{p{\linewidth}}}]{
                \texttt{\tiny{[P1$\langle$GM]}}
                \texttt{Ceza: (1,6) boş değil; "─" içeriyor.} \\
\\ 
\texttt{Toplamda 16 cezanın 14 kadarını biriktirdiniz.Lütfen tekrar deneyin!} \\
            }
        }
    }
    & \\ \\

    \theutterance \stepcounter{utterance}  
    & \multicolumn{4}{p{0.6\linewidth}}{
        \cellcolor[rgb]{0.9,0.9,0.9}{
            \makecell[{{p{\linewidth}}}]{
                \texttt{\tiny{[P1$\rangle$GM]}}
                \texttt{SÖYLE: (1,6) da ─ var, boş değil. Benim tahtamda 1. satırda boş koordinatlar: (1,1), (2,1) ve 7. satırda (3,7), (4,7), (5,7). S için (3,7) veya M için (3,7) düşünelim. S ve M için hangi koordinatlar sende boş? Belki S ve M'yi 7. satıra taşıyalım: S(3,7), M(5,7) gibi.} \\
            }
        }
    }
    & & \\ \\

    \theutterance \stepcounter{utterance}  
    & & & \multicolumn{2}{p{0.3\linewidth}}{
        \cellcolor[rgb]{0.9,0.9,0.9}{
            \makecell[{{p{\linewidth}}}]{
                \texttt{\tiny{[GM$|$GM]}}
                \texttt{SÖYLE: (1,6) da ─ var, boş değil. Benim tahtamda 1. satırda boş koordinatlar: (1,1), (2,1) ve 7. satırda (3,7), (4,7), (5,7). S için (3,7) veya M için (3,7) düşünelim. S ve M için hangi koordinatlar sende boş? Belki S ve M'yi 7. satıra taşıyalım: S(3,7), M(5,7) gibi.} \\
            }
        }
    }
    & & \\ \\

    \theutterance \stepcounter{utterance}  
    & & \multicolumn{4}{p{0.6\linewidth}}{
        \cellcolor[rgb]{0.9,0.9,0.9}{
            \makecell[{{p{\linewidth}}}]{
                \texttt{\tiny{[GM$\rangle$P2]}}
                \texttt{Mesajın diğer oyuncuya iletildi.} \\
\\ 
\\ 
\texttt{Şu anda en fazla 28 turdan 22. turu oynuyorsunuz.} \\
\texttt{Toplamda 16 cezanın 14 kadarını biriktirdiniz.Diğer oyuncu şu mesajı gönderdi:} \\
\\ 
\\ 
\texttt{"(1,6) da ─ var, boş değil. Benim tahtamda 1. satırda boş koordinatlar: (1,1), (2,1) ve 7. satırda (3,7), (4,7), (5,7). S için (3,7) veya M için (3,7) düşünelim. S ve M için hangi koordinatlar sende boş? Belki S ve M'yi 7. satıra taşıyalım: S(3,7), M(5,7) gibi."} \\
\\ 
\\ 
\texttt{Bir sonraki komutun nedir?} \\
            }
        }
    }
    & \\ \\

    \theutterance \stepcounter{utterance}  
    & & & \multicolumn{4}{p{0.6\linewidth}}{
        \cellcolor[rgb]{0.9,0.9,0.9}{
            \makecell[{{p{\linewidth}}}]{
                \texttt{\tiny{[GM$\langle$P2]}}
                \texttt{SÖYLE: (3,7) ve (5,7) benim tahtamda boş. S için (3,7), M için (5,7) uygun. Güncel hedefler: P(5,4), T(7,3), D(2,3), A(1,3), S(3,7), M(5,7). S'yi (3,7)'ye taşıyorum.} \\
            }
        }
    }
     \\ \\

    \theutterance \stepcounter{utterance}  
    & & & \multicolumn{2}{p{0.3\linewidth}}{
        \cellcolor[rgb]{0.9,0.9,0.9}{
            \makecell[{{p{\linewidth}}}]{
                \texttt{\tiny{[GM$|$GM]}}
                \texttt{SÖYLE: (3,7) ve (5,7) benim tahtamda boş. S için (3,7), M için (5,7) uygun. Güncel hedefler: P(5,4), T(7,3), D(2,3), A(1,3), S(3,7), M(5,7). S'yi (3,7)'ye taşıyorum.} \\
            }
        }
    }
    & & \\ \\

    \theutterance \stepcounter{utterance}  
    & & \multicolumn{4}{p{0.6\linewidth}}{
        \cellcolor[rgb]{0.9,0.9,0.9}{
            \makecell[{{p{\linewidth}}}]{
                \texttt{\tiny{[P1$\langle$GM]}}
                \texttt{Mesajın diğer oyuncuya iletildi.} \\
\\ 
\\ 
\texttt{Şu anda en fazla 28 turdan 22. turu oynuyorsunuz.} \\
\texttt{Toplamda 16 cezanın 14 kadarını biriktirdiniz.Diğer oyuncu şu mesajı gönderdi:} \\
\\ 
\\ 
\texttt{"(3,7) ve (5,7) benim tahtamda boş. S için (3,7), M için (5,7) uygun. Güncel hedefler: P(5,4), T(7,3), D(2,3), A(1,3), S(3,7), M(5,7). S'yi (3,7)'ye taşıyorum."} \\
\\ 
\\ 
\texttt{Bir sonraki komutun nedir?} \\
            }
        }
    }
    & \\ \\

    \theutterance \stepcounter{utterance}  
    & \multicolumn{4}{p{0.6\linewidth}}{
        \cellcolor[rgb]{0.9,0.9,0.9}{
            \makecell[{{p{\linewidth}}}]{
                \texttt{\tiny{[P1$\rangle$GM]}}
                \texttt{HAMLE: S, (3,7)} \\
            }
        }
    }
    & & \\ \\

    \theutterance \stepcounter{utterance}  
    & & & \multicolumn{2}{p{0.3\linewidth}}{
        \cellcolor[rgb]{0.9,0.9,0.9}{
            \makecell[{{p{\linewidth}}}]{
                \texttt{\tiny{[GM$|$GM]}}
                \texttt{HAMLE: S, (3,7)} \\
            }
        }
    }
    & & \\ \\

    \theutterance \stepcounter{utterance}  
    & & & \multicolumn{2}{p{0.3\linewidth}}{
        \cellcolor[rgb]{0.9,0.9,0.9}{
            \makecell[{{p{\linewidth}}}]{
                \texttt{\tiny{[GM$|$GM]}}
                \texttt{Ceza: (3,7) boş değil; "S" içeriyor.} \\
\\ 
\texttt{Toplamda 16 cezanın 15 kadarını biriktirdiniz.Lütfen tekrar deneyin!} \\
            }
        }
    }
    & & \\ \\

    \theutterance \stepcounter{utterance}  
    & & \multicolumn{4}{p{0.6\linewidth}}{
        \cellcolor[rgb]{0.9,0.9,0.9}{
            \makecell[{{p{\linewidth}}}]{
                \texttt{\tiny{[P1$\langle$GM]}}
                \texttt{Ceza: (3,7) boş değil; "S" içeriyor.} \\
\\ 
\texttt{Toplamda 16 cezanın 15 kadarını biriktirdiniz.Lütfen tekrar deneyin!} \\
            }
        }
    }
    & \\ \\

    \theutterance \stepcounter{utterance}  
    & \multicolumn{4}{p{0.6\linewidth}}{
        \cellcolor[rgb]{0.9,0.9,0.9}{
            \makecell[{{p{\linewidth}}}]{
                \texttt{\tiny{[P1$\rangle$GM]}}
                \texttt{SÖYLE: (3,7) de S'nin şu anki konumu zaten! Benim tahtamda S (2,7)'de. S'yi taşımak istiyorum. (3,7) boş değil çünkü orada S yazıyor {-} pardon, S (2,7)'de. Bir daha kontrol edeyim: Benim tahtamda 7. satır: (1,7) boş, (2,7) S var, (3,7) boş mu? Izgaraya bakıyorum: 7. satırda ║◌│S│◌│◌║ yani (3,7) boş görünüyor, (4,7) de boş. Ama az önce boş değil dedi. (3,7)'de ne var tam olarak?} \\
            }
        }
    }
    & & \\ \\

    \theutterance \stepcounter{utterance}  
    & & & \multicolumn{2}{p{0.3\linewidth}}{
        \cellcolor[rgb]{0.9,0.9,0.9}{
            \makecell[{{p{\linewidth}}}]{
                \texttt{\tiny{[GM$|$GM]}}
                \texttt{SÖYLE: (3,7) de S'nin şu anki konumu zaten! Benim tahtamda S (2,7)'de. S'yi taşımak istiyorum. (3,7) boş değil çünkü orada S yazıyor {-} pardon, S (2,7)'de. Bir daha kontrol edeyim: Benim tahtamda 7. satır: (1,7) boş, (2,7) S var, (3,7) boş mu? Izgaraya bakıyorum: 7. satırda ║◌│S│◌│◌║ yani (3,7) boş görünüyor, (4,7) de boş. Ama az önce boş değil dedi. (3,7)'de ne var tam olarak?} \\
            }
        }
    }
    & & \\ \\

    \theutterance \stepcounter{utterance}  
    & & \multicolumn{4}{p{0.6\linewidth}}{
        \cellcolor[rgb]{0.9,0.9,0.9}{
            \makecell[{{p{\linewidth}}}]{
                \texttt{\tiny{[GM$\rangle$P2]}}
                \texttt{Mesajın diğer oyuncuya iletildi.} \\
\\ 
\\ 
\texttt{Şu anda en fazla 28 turdan 23. turu oynuyorsunuz.} \\
\texttt{Toplamda 16 cezanın 15 kadarını biriktirdiniz.Diğer oyuncu şu mesajı gönderdi:} \\
\\ 
\\ 
\texttt{"(3,7) de S'nin şu anki konumu zaten! Benim tahtamda S (2,7)'de. S'yi taşımak istiyorum. (3,7) boş değil çünkü orada S yazıyor {-} pardon, S (2,7)'de. Bir daha kontrol edeyim: Benim tahtamda 7. satır: (1,7) boş, (2,7) S var, (3,7) boş mu? Izgaraya bakıyorum: 7. satırda ║◌│S│◌│◌║ yani (3,7) boş görünüyor, (4,7) de boş. Ama az önce boş değil dedi. (3,7)'de ne var tam olarak?"} \\
\\ 
\\ 
\texttt{Bir sonraki komutun nedir?} \\
            }
        }
    }
    & \\ \\

    \theutterance \stepcounter{utterance}  
    & & & \multicolumn{4}{p{0.6\linewidth}}{
        \cellcolor[rgb]{0.9,0.9,0.9}{
            \makecell[{{p{\linewidth}}}]{
                \texttt{\tiny{[GM$\langle$P2]}}
                \texttt{SÖYLE: Benim tahtamda 7. satır: M(1,7), (2,7) boş, (3,7) boş, (4,7) boş, (5,7) boş, (6,7) boş, (7,7) boş (T'yi taşıdım). S için (3,7) bende boş. Sende de boşsa S'yi (3,7)'ye taşıyalım. M için (5,7) ikimizde de boş. Devam edelim, cezalar birikiyor.} \\
            }
        }
    }
     \\ \\

    \theutterance \stepcounter{utterance}  
    & & & \multicolumn{2}{p{0.3\linewidth}}{
        \cellcolor[rgb]{0.9,0.9,0.9}{
            \makecell[{{p{\linewidth}}}]{
                \texttt{\tiny{[GM$|$GM]}}
                \texttt{SÖYLE: Benim tahtamda 7. satır: M(1,7), (2,7) boş, (3,7) boş, (4,7) boş, (5,7) boş, (6,7) boş, (7,7) boş (T'yi taşıdım). S için (3,7) bende boş. Sende de boşsa S'yi (3,7)'ye taşıyalım. M için (5,7) ikimizde de boş. Devam edelim, cezalar birikiyor.} \\
            }
        }
    }
    & & \\ \\

    \theutterance \stepcounter{utterance}  
    & & \multicolumn{4}{p{0.6\linewidth}}{
        \cellcolor[rgb]{0.9,0.9,0.9}{
            \makecell[{{p{\linewidth}}}]{
                \texttt{\tiny{[P1$\langle$GM]}}
                \texttt{Mesajın diğer oyuncuya iletildi.} \\
\\ 
\\ 
\texttt{Şu anda en fazla 28 turdan 23. turu oynuyorsunuz.} \\
\texttt{Toplamda 16 cezanın 15 kadarını biriktirdiniz.Diğer oyuncu şu mesajı gönderdi:} \\
\\ 
\\ 
\texttt{"Benim tahtamda 7. satır: M(1,7), (2,7) boş, (3,7) boş, (4,7) boş, (5,7) boş, (6,7) boş, (7,7) boş (T'yi taşıdım). S için (3,7) bende boş. Sende de boşsa S'yi (3,7)'ye taşıyalım. M için (5,7) ikimizde de boş. Devam edelim, cezalar birikiyor."} \\
\\ 
\\ 
\texttt{Bir sonraki komutun nedir?} \\
            }
        }
    }
    & \\ \\

    \theutterance \stepcounter{utterance}  
    & \multicolumn{4}{p{0.6\linewidth}}{
        \cellcolor[rgb]{0.9,0.9,0.9}{
            \makecell[{{p{\linewidth}}}]{
                \texttt{\tiny{[P1$\rangle$GM]}}
                \texttt{HAMLE: S, (3,7)} \\
            }
        }
    }
    & & \\ \\

    \theutterance \stepcounter{utterance}  
    & & & \multicolumn{2}{p{0.3\linewidth}}{
        \cellcolor[rgb]{0.9,0.9,0.9}{
            \makecell[{{p{\linewidth}}}]{
                \texttt{\tiny{[GM$|$GM]}}
                \texttt{HAMLE: S, (3,7)} \\
            }
        }
    }
    & & \\ \\

    \theutterance \stepcounter{utterance}  
    & & & \multicolumn{2}{p{0.3\linewidth}}{
        \cellcolor[rgb]{0.9,0.9,0.9}{
            \makecell[{{p{\linewidth}}}]{
                \texttt{\tiny{[GM$|$GM]}}
                \texttt{Ceza: (3,7) boş değil; "S" içeriyor.} \\
\\ 
\texttt{Toplamda 16 cezanın 16 kadarını biriktirdiniz.Lütfen tekrar deneyin!} \\
            }
        }
    }
    & & \\ \\

    \theutterance \stepcounter{utterance}  
    & & \multicolumn{4}{p{0.6\linewidth}}{
        \cellcolor[rgb]{0.9,0.9,0.9}{
            \makecell[{{p{\linewidth}}}]{
                \texttt{\tiny{[P1$\langle$GM]}}
                \texttt{Ceza: (3,7) boş değil; "S" içeriyor.} \\
\\ 
\texttt{Toplamda 16 cezanın 16 kadarını biriktirdiniz.Lütfen tekrar deneyin!} \\
            }
        }
    }
    & \\ \\

    \theutterance \stepcounter{utterance}  
    & \multicolumn{4}{p{0.6\linewidth}}{
        \cellcolor[rgb]{0.9,0.9,0.9}{
            \makecell[{{p{\linewidth}}}]{
                \texttt{\tiny{[P1$\rangle$GM]}}
                \texttt{SÖYLE: Tamam sanırım hata yapıyorum. Benim tahtamda S şu an (2,7)'de. (3,7)'ye taşımaya çalışıyorum ama sistem boş değil diyor. Oysa güncel ızgaramda (3,7) boş görünüyor. S'yi önce (2,7)'den başka bir yere taşıyıp sonra (3,7)'ye mi taşımam gerekiyor? Ya da direkt (3,7) boşsa neden taşıyamıyorum, kontrol edebilir misin?} \\
            }
        }
    }
    & & \\ \\

    \theutterance \stepcounter{utterance}  
    & & & \multicolumn{2}{p{0.3\linewidth}}{
        \cellcolor[rgb]{0.9,0.9,0.9}{
            \makecell[{{p{\linewidth}}}]{
                \texttt{\tiny{[GM$|$GM]}}
                \texttt{SÖYLE: Tamam sanırım hata yapıyorum. Benim tahtamda S şu an (2,7)'de. (3,7)'ye taşımaya çalışıyorum ama sistem boş değil diyor. Oysa güncel ızgaramda (3,7) boş görünüyor. S'yi önce (2,7)'den başka bir yere taşıyıp sonra (3,7)'ye mi taşımam gerekiyor? Ya da direkt (3,7) boşsa neden taşıyamıyorum, kontrol edebilir misin?} \\
            }
        }
    }
    & & \\ \\

    \theutterance \stepcounter{utterance}  
    & & \multicolumn{4}{p{0.6\linewidth}}{
        \cellcolor[rgb]{0.9,0.9,0.9}{
            \makecell[{{p{\linewidth}}}]{
                \texttt{\tiny{[GM$\rangle$P2]}}
                \texttt{Mesajın diğer oyuncuya iletildi.} \\
\\ 
\\ 
\texttt{Şu anda en fazla 28 turdan 24. turu oynuyorsunuz.} \\
\texttt{Toplamda 16 cezanın 16 kadarını biriktirdiniz.Diğer oyuncu şu mesajı gönderdi:} \\
\\ 
\\ 
\texttt{"Tamam sanırım hata yapıyorum. Benim tahtamda S şu an (2,7)'de. (3,7)'ye taşımaya çalışıyorum ama sistem boş değil diyor. Oysa güncel ızgaramda (3,7) boş görünüyor. S'yi önce (2,7)'den başka bir yere taşıyıp sonra (3,7)'ye mi taşımam gerekiyor? Ya da direkt (3,7) boşsa neden taşıyamıyorum, kontrol edebilir misin?"} \\
\\ 
\\ 
\texttt{Bir sonraki komutun nedir?} \\
            }
        }
    }
    & \\ \\

    \theutterance \stepcounter{utterance}  
    & & & \multicolumn{4}{p{0.6\linewidth}}{
        \cellcolor[rgb]{0.9,0.9,0.9}{
            \makecell[{{p{\linewidth}}}]{
                \texttt{\tiny{[GM$\langle$P2]}}
                \texttt{SÖYLE: 16 cezaya ulaştık, oyunu kaybetme riskimiz var. Hemen toparlayalım. S'yi (3,7) yerine (5,7) dene, M'yi de (3,7) yapalım. Veya direkt olarak S'yi (1,6) dene, M'yi (5,6) dene. Nerede boş olduğunu sen daha iyi görüyorsun. S ve M için sende kesin boş olan iki koordinat söyle, ben de uyumlu olanı seçeyim. Hızlı olalım.} \\
            }
        }
    }
     \\ \\

    \theutterance \stepcounter{utterance}  
    & & & \multicolumn{2}{p{0.3\linewidth}}{
        \cellcolor[rgb]{0.9,0.9,0.9}{
            \makecell[{{p{\linewidth}}}]{
                \texttt{\tiny{[GM$|$GM]}}
                \texttt{SÖYLE: 16 cezaya ulaştık, oyunu kaybetme riskimiz var. Hemen toparlayalım. S'yi (3,7) yerine (5,7) dene, M'yi de (3,7) yapalım. Veya direkt olarak S'yi (1,6) dene, M'yi (5,6) dene. Nerede boş olduğunu sen daha iyi görüyorsun. S ve M için sende kesin boş olan iki koordinat söyle, ben de uyumlu olanı seçeyim. Hızlı olalım.} \\
            }
        }
    }
    & & \\ \\

    \theutterance \stepcounter{utterance}  
    & & \multicolumn{4}{p{0.6\linewidth}}{
        \cellcolor[rgb]{0.9,0.9,0.9}{
            \makecell[{{p{\linewidth}}}]{
                \texttt{\tiny{[P1$\langle$GM]}}
                \texttt{Mesajın diğer oyuncuya iletildi.} \\
\\ 
\\ 
\texttt{Şu anda en fazla 28 turdan 24. turu oynuyorsunuz.} \\
\texttt{Toplamda 16 cezanın 16 kadarını biriktirdiniz.Diğer oyuncu şu mesajı gönderdi:} \\
\\ 
\\ 
\texttt{"16 cezaya ulaştık, oyunu kaybetme riskimiz var. Hemen toparlayalım. S'yi (3,7) yerine (5,7) dene, M'yi de (3,7) yapalım. Veya direkt olarak S'yi (1,6) dene, M'yi (5,6) dene. Nerede boş olduğunu sen daha iyi görüyorsun. S ve M için sende kesin boş olan iki koordinat söyle, ben de uyumlu olanı seçeyim. Hızlı olalım."} \\
\\ 
\\ 
\texttt{Bir sonraki komutun nedir?} \\
            }
        }
    }
    & \\ \\

    \theutterance \stepcounter{utterance}  
    & \multicolumn{4}{p{0.6\linewidth}}{
        \cellcolor[rgb]{0.9,0.9,0.9}{
            \makecell[{{p{\linewidth}}}]{
                \texttt{\tiny{[P1$\rangle$GM]}}
                \texttt{SÖYLE: Izgarama tekrar bakıyorum: (1,1) boş, (2,1) boş, (4,1) boş, (5,1) boş, (6,1) boş, (7,1) boş, (1,2) boş, (5,2) boş, (7,2) boş, (1,4) boş, (5,5) boş, (6,5) boş, (7,5) boş, (1,6) ─ var, (2,6) boş, (3,7) boş görünüyor ama hata verdi, (4,7) boş, (5,7) boş, (6,7) boş, (7,7) boş. S ve M için (4,7) ve (5,7) öneriyorum. S(4,7), M(5,7) olsun.} \\
            }
        }
    }
    & & \\ \\

    \theutterance \stepcounter{utterance}  
    & & & \multicolumn{2}{p{0.3\linewidth}}{
        \cellcolor[rgb]{0.9,0.9,0.9}{
            \makecell[{{p{\linewidth}}}]{
                \texttt{\tiny{[GM$|$GM]}}
                \texttt{SÖYLE: Izgarama tekrar bakıyorum: (1,1) boş, (2,1) boş, (4,1) boş, (5,1) boş, (6,1) boş, (7,1) boş, (1,2) boş, (5,2) boş, (7,2) boş, (1,4) boş, (5,5) boş, (6,5) boş, (7,5) boş, (1,6) ─ var, (2,6) boş, (3,7) boş görünüyor ama hata verdi, (4,7) boş, (5,7) boş, (6,7) boş, (7,7) boş. S ve M için (4,7) ve (5,7) öneriyorum. S(4,7), M(5,7) olsun.} \\
            }
        }
    }
    & & \\ \\

    \theutterance \stepcounter{utterance}  
    & & \multicolumn{4}{p{0.6\linewidth}}{
        \cellcolor[rgb]{0.9,0.9,0.9}{
            \makecell[{{p{\linewidth}}}]{
                \texttt{\tiny{[GM$\rangle$P2]}}
                \texttt{Mesajın diğer oyuncuya iletildi.} \\
\\ 
\\ 
\texttt{Şu anda en fazla 28 turdan 25. turu oynuyorsunuz.} \\
\texttt{Toplamda 16 cezanın 16 kadarını biriktirdiniz.Diğer oyuncu şu mesajı gönderdi:} \\
\\ 
\\ 
\texttt{"Izgarama tekrar bakıyorum: (1,1) boş, (2,1) boş, (4,1) boş, (5,1) boş, (6,1) boş, (7,1) boş, (1,2) boş, (5,2) boş, (7,2) boş, (1,4) boş, (5,5) boş, (6,5) boş, (7,5) boş, (1,6) ─ var, (2,6) boş, (3,7) boş görünüyor ama hata verdi, (4,7) boş, (5,7) boş, (6,7) boş, (7,7) boş. S ve M için (4,7) ve (5,7) öneriyorum. S(4,7), M(5,7) olsun."} \\
\\ 
\\ 
\texttt{Bir sonraki komutun nedir?} \\
            }
        }
    }
    & \\ \\

    \theutterance \stepcounter{utterance}  
    & & & \multicolumn{4}{p{0.6\linewidth}}{
        \cellcolor[rgb]{0.9,0.9,0.9}{
            \makecell[{{p{\linewidth}}}]{
                \texttt{\tiny{[GM$\langle$P2]}}
                \texttt{SÖYLE: (4,7) ve (5,7) benim tahtamda boş. S(4,7), M(5,7) kabul. S'yi (4,7)'ye taşı. Ben de S'yi (4,7)'ye taşıyorum.} \\
            }
        }
    }
     \\ \\

    \theutterance \stepcounter{utterance}  
    & & & \multicolumn{2}{p{0.3\linewidth}}{
        \cellcolor[rgb]{0.9,0.9,0.9}{
            \makecell[{{p{\linewidth}}}]{
                \texttt{\tiny{[GM$|$GM]}}
                \texttt{SÖYLE: (4,7) ve (5,7) benim tahtamda boş. S(4,7), M(5,7) kabul. S'yi (4,7)'ye taşı. Ben de S'yi (4,7)'ye taşıyorum.} \\
            }
        }
    }
    & & \\ \\

    \theutterance \stepcounter{utterance}  
    & & \multicolumn{4}{p{0.6\linewidth}}{
        \cellcolor[rgb]{0.9,0.9,0.9}{
            \makecell[{{p{\linewidth}}}]{
                \texttt{\tiny{[P1$\langle$GM]}}
                \texttt{Mesajın diğer oyuncuya iletildi.} \\
\\ 
\\ 
\texttt{Şu anda en fazla 28 turdan 25. turu oynuyorsunuz.} \\
\texttt{Toplamda 16 cezanın 16 kadarını biriktirdiniz.Diğer oyuncu şu mesajı gönderdi:} \\
\\ 
\\ 
\texttt{"(4,7) ve (5,7) benim tahtamda boş. S(4,7), M(5,7) kabul. S'yi (4,7)'ye taşı. Ben de S'yi (4,7)'ye taşıyorum."} \\
\\ 
\\ 
\texttt{Bir sonraki komutun nedir?} \\
            }
        }
    }
    & \\ \\

    \theutterance \stepcounter{utterance}  
    & \multicolumn{4}{p{0.6\linewidth}}{
        \cellcolor[rgb]{0.9,0.9,0.9}{
            \makecell[{{p{\linewidth}}}]{
                \texttt{\tiny{[P1$\rangle$GM]}}
                \texttt{HAMLE: S, (4,7)} \\
            }
        }
    }
    & & \\ \\

    \theutterance \stepcounter{utterance}  
    & & & \multicolumn{2}{p{0.3\linewidth}}{
        \cellcolor[rgb]{0.9,0.9,0.9}{
            \makecell[{{p{\linewidth}}}]{
                \texttt{\tiny{[GM$|$GM]}}
                \texttt{HAMLE: S, (4,7)} \\
            }
        }
    }
    & & \\ \\

    \theutterance \stepcounter{utterance}  
    & & & \multicolumn{2}{p{0.3\linewidth}}{
        \cellcolor[rgb]{0.9,0.9,0.9}{
            \makecell[{{p{\linewidth}}}]{
                \texttt{\tiny{[GM$|$GM]}}
                \texttt{Ceza: (4,7) boş değil; "│" içeriyor.} \\
\\ 
\texttt{Toplamda 16 cezanın 17 kadarını biriktirdiniz.Lütfen tekrar deneyin!} \\
            }
        }
    }
    & & \\ \\

    \theutterance \stepcounter{utterance}  
    & & \multicolumn{4}{p{0.6\linewidth}}{
        \cellcolor[rgb]{0.9,0.9,0.9}{
            \makecell[{{p{\linewidth}}}]{
                \texttt{\tiny{[P1$\langle$GM]}}
                \texttt{Ceza: (4,7) boş değil; "│" içeriyor.} \\
\\ 
\texttt{Toplamda 16 cezanın 17 kadarını biriktirdiniz.Lütfen tekrar deneyin!} \\
            }
        }
    }
    & \\ \\

    \theutterance \stepcounter{utterance}  
    & \multicolumn{4}{p{0.6\linewidth}}{
        \cellcolor[rgb]{0.9,0.9,0.9}{
            \makecell[{{p{\linewidth}}}]{
                \texttt{\tiny{[P1$\rangle$GM]}}
                \texttt{SÖYLE: (4,7) de boş değilmiş, │ var. 7. satırda (4,7) ve (5,7) dışında (6,7) ve (7,7) boş. Ayrıca (2,6) ve (5,5) de boş. S ve M için (6,7) ve (7,7) uygun mu? Veya (2,6) ve (5,5)? Son 3 turumuz kaldı, hızlı karar verelim.} \\
            }
        }
    }
    & & \\ \\

    \theutterance \stepcounter{utterance}  
    & & & \multicolumn{2}{p{0.3\linewidth}}{
        \cellcolor[rgb]{0.9,0.9,0.9}{
            \makecell[{{p{\linewidth}}}]{
                \texttt{\tiny{[GM$|$GM]}}
                \texttt{SÖYLE: (4,7) de boş değilmiş, │ var. 7. satırda (4,7) ve (5,7) dışında (6,7) ve (7,7) boş. Ayrıca (2,6) ve (5,5) de boş. S ve M için (6,7) ve (7,7) uygun mu? Veya (2,6) ve (5,5)? Son 3 turumuz kaldı, hızlı karar verelim.} \\
            }
        }
    }
    & & \\ \\

    \theutterance \stepcounter{utterance}  
    & & \multicolumn{4}{p{0.6\linewidth}}{
        \cellcolor[rgb]{0.9,0.9,0.9}{
            \makecell[{{p{\linewidth}}}]{
                \texttt{\tiny{[GM$\rangle$P2]}}
                \texttt{Mesajın diğer oyuncuya iletildi.} \\
\\ 
\\ 
\texttt{Şu anda en fazla 28 turdan 26. turu oynuyorsunuz.} \\
\texttt{Toplamda 16 cezanın 17 kadarını biriktirdiniz.Diğer oyuncu şu mesajı gönderdi:} \\
\\ 
\\ 
\texttt{"(4,7) de boş değilmiş, │ var. 7. satırda (4,7) ve (5,7) dışında (6,7) ve (7,7) boş. Ayrıca (2,6) ve (5,5) de boş. S ve M için (6,7) ve (7,7) uygun mu? Veya (2,6) ve (5,5)? Son 3 turumuz kaldı, hızlı karar verelim."} \\
\\ 
\\ 
\texttt{Bir sonraki komutun nedir?} \\
            }
        }
    }
    & \\ \\

    \theutterance \stepcounter{utterance}  
    & & & \multicolumn{4}{p{0.6\linewidth}}{
        \cellcolor[rgb]{0.9,0.9,0.9}{
            \makecell[{{p{\linewidth}}}]{
                \texttt{\tiny{[GM$\langle$P2]}}
                \texttt{SÖYLE: (6,7) ve (7,7) benim tahtamda boş. S(6,7), M(7,7) yapalım. Hemen taşıyalım.} \\
            }
        }
    }
     \\ \\

    \theutterance \stepcounter{utterance}  
    & & & \multicolumn{2}{p{0.3\linewidth}}{
        \cellcolor[rgb]{0.9,0.9,0.9}{
            \makecell[{{p{\linewidth}}}]{
                \texttt{\tiny{[GM$|$GM]}}
                \texttt{SÖYLE: (6,7) ve (7,7) benim tahtamda boş. S(6,7), M(7,7) yapalım. Hemen taşıyalım.} \\
            }
        }
    }
    & & \\ \\

    \theutterance \stepcounter{utterance}  
    & & \multicolumn{4}{p{0.6\linewidth}}{
        \cellcolor[rgb]{0.9,0.9,0.9}{
            \makecell[{{p{\linewidth}}}]{
                \texttt{\tiny{[P1$\langle$GM]}}
                \texttt{Mesajın diğer oyuncuya iletildi.} \\
\\ 
\\ 
\texttt{Şu anda en fazla 28 turdan 26. turu oynuyorsunuz.} \\
\texttt{Toplamda 16 cezanın 17 kadarını biriktirdiniz.Diğer oyuncu şu mesajı gönderdi:} \\
\\ 
\\ 
\texttt{"(6,7) ve (7,7) benim tahtamda boş. S(6,7), M(7,7) yapalım. Hemen taşıyalım."} \\
\\ 
\\ 
\texttt{Bir sonraki komutun nedir?} \\
            }
        }
    }
    & \\ \\

    \theutterance \stepcounter{utterance}  
    & \multicolumn{4}{p{0.6\linewidth}}{
        \cellcolor[rgb]{0.9,0.9,0.9}{
            \makecell[{{p{\linewidth}}}]{
                \texttt{\tiny{[P1$\rangle$GM]}}
                \texttt{HAMLE: S, (6,7)} \\
            }
        }
    }
    & & \\ \\

    \theutterance \stepcounter{utterance}  
    & & & \multicolumn{2}{p{0.3\linewidth}}{
        \cellcolor[rgb]{0.9,0.9,0.9}{
            \makecell[{{p{\linewidth}}}]{
                \texttt{\tiny{[GM$|$GM]}}
                \texttt{HAMLE: S, (6,7)} \\
            }
        }
    }
    & & \\ \\

    \theutterance \stepcounter{utterance}  
    & & & \multicolumn{2}{p{0.3\linewidth}}{
        \cellcolor[rgb]{0.9,0.9,0.9}{
            \makecell[{{p{\linewidth}}}]{
                \texttt{\tiny{[GM$|$GM]}}
                \texttt{Ceza: (6,7) boş değil; "│" içeriyor.} \\
\\ 
\texttt{Toplamda 16 cezanın 18 kadarını biriktirdiniz.Lütfen tekrar deneyin!} \\
            }
        }
    }
    & & \\ \\

    \theutterance \stepcounter{utterance}  
    & & \multicolumn{4}{p{0.6\linewidth}}{
        \cellcolor[rgb]{0.9,0.9,0.9}{
            \makecell[{{p{\linewidth}}}]{
                \texttt{\tiny{[P1$\langle$GM]}}
                \texttt{Ceza: (6,7) boş değil; "│" içeriyor.} \\
\\ 
\texttt{Toplamda 16 cezanın 18 kadarını biriktirdiniz.Lütfen tekrar deneyin!} \\
            }
        }
    }
    & \\ \\

    \theutterance \stepcounter{utterance}  
    & \multicolumn{4}{p{0.6\linewidth}}{
        \cellcolor[rgb]{0.9,0.9,0.9}{
            \makecell[{{p{\linewidth}}}]{
                \texttt{\tiny{[P1$\rangle$GM]}}
                \texttt{SÖYLE: 7. satırın tamamı boş değil galiba. (1,7) boş, (5,5) boş, (2,6) boş, (5,2) boş. S ve M için (1,7) ve (5,5) ya da (2,6) ve (5,2) kullanalım. Son 2 tur! S(1,7), M(5,5) diyorum, sende boş mu?} \\
            }
        }
    }
    & & \\ \\

    \theutterance \stepcounter{utterance}  
    & & & \multicolumn{2}{p{0.3\linewidth}}{
        \cellcolor[rgb]{0.9,0.9,0.9}{
            \makecell[{{p{\linewidth}}}]{
                \texttt{\tiny{[GM$|$GM]}}
                \texttt{SÖYLE: 7. satırın tamamı boş değil galiba. (1,7) boş, (5,5) boş, (2,6) boş, (5,2) boş. S ve M için (1,7) ve (5,5) ya da (2,6) ve (5,2) kullanalım. Son 2 tur! S(1,7), M(5,5) diyorum, sende boş mu?} \\
            }
        }
    }
    & & \\ \\

    \theutterance \stepcounter{utterance}  
    & & \multicolumn{4}{p{0.6\linewidth}}{
        \cellcolor[rgb]{0.9,0.9,0.9}{
            \makecell[{{p{\linewidth}}}]{
                \texttt{\tiny{[GM$\rangle$P2]}}
                \texttt{Mesajın diğer oyuncuya iletildi.} \\
\\ 
\\ 
\texttt{Şu anda en fazla 28 turdan 27. turu oynuyorsunuz.} \\
\texttt{Toplamda 16 cezanın 18 kadarını biriktirdiniz.Diğer oyuncu şu mesajı gönderdi:} \\
\\ 
\\ 
\texttt{"7. satırın tamamı boş değil galiba. (1,7) boş, (5,5) boş, (2,6) boş, (5,2) boş. S ve M için (1,7) ve (5,5) ya da (2,6) ve (5,2) kullanalım. Son 2 tur! S(1,7), M(5,5) diyorum, sende boş mu?"} \\
\\ 
\\ 
\texttt{Bir sonraki komutun nedir?} \\
            }
        }
    }
    & \\ \\

    \theutterance \stepcounter{utterance}  
    & & & \multicolumn{4}{p{0.6\linewidth}}{
        \cellcolor[rgb]{0.9,0.9,0.9}{
            \makecell[{{p{\linewidth}}}]{
                \texttt{\tiny{[GM$\langle$P2]}}
                \texttt{SÖYLE: (1,7) benim tahtamda M var, boş değil. (5,5) boş. (2,6) boş, (5,2) boş. S(2,6) ve M(5,5) yapalım. Veya S(5,2) ve M(2,6). Sen seç, hemen söyle.} \\
            }
        }
    }
     \\ \\

    \theutterance \stepcounter{utterance}  
    & & & \multicolumn{2}{p{0.3\linewidth}}{
        \cellcolor[rgb]{0.9,0.9,0.9}{
            \makecell[{{p{\linewidth}}}]{
                \texttt{\tiny{[GM$|$GM]}}
                \texttt{SÖYLE: (1,7) benim tahtamda M var, boş değil. (5,5) boş. (2,6) boş, (5,2) boş. S(2,6) ve M(5,5) yapalım. Veya S(5,2) ve M(2,6). Sen seç, hemen söyle.} \\
            }
        }
    }
    & & \\ \\

    \theutterance \stepcounter{utterance}  
    & & \multicolumn{4}{p{0.6\linewidth}}{
        \cellcolor[rgb]{0.9,0.9,0.9}{
            \makecell[{{p{\linewidth}}}]{
                \texttt{\tiny{[P1$\langle$GM]}}
                \texttt{Mesajın diğer oyuncuya iletildi.} \\
\\ 
\\ 
\texttt{Şu anda en fazla 28 turdan 27. turu oynuyorsunuz.} \\
\texttt{Toplamda 16 cezanın 18 kadarını biriktirdiniz.Diğer oyuncu şu mesajı gönderdi:} \\
\\ 
\\ 
\texttt{"(1,7) benim tahtamda M var, boş değil. (5,5) boş. (2,6) boş, (5,2) boş. S(2,6) ve M(5,5) yapalım. Veya S(5,2) ve M(2,6). Sen seç, hemen söyle."} \\
\\ 
\\ 
\texttt{Bir sonraki komutun nedir?} \\
            }
        }
    }
    & \\ \\

    \theutterance \stepcounter{utterance}  
    & \multicolumn{4}{p{0.6\linewidth}}{
        \cellcolor[rgb]{0.9,0.9,0.9}{
            \makecell[{{p{\linewidth}}}]{
                \texttt{\tiny{[P1$\rangle$GM]}}
                \texttt{SÖYLE: S(2,6) ve M(5,5) olsun. Hemen dene!} \\
            }
        }
    }
    & & \\ \\

    \theutterance \stepcounter{utterance}  
    & & & \multicolumn{2}{p{0.3\linewidth}}{
        \cellcolor[rgb]{0.9,0.9,0.9}{
            \makecell[{{p{\linewidth}}}]{
                \texttt{\tiny{[GM$|$GM]}}
                \texttt{SÖYLE: S(2,6) ve M(5,5) olsun. Hemen dene!} \\
            }
        }
    }
    & & \\ \\

    \theutterance \stepcounter{utterance}  
    & & \multicolumn{4}{p{0.6\linewidth}}{
        \cellcolor[rgb]{0.9,0.9,0.9}{
            \makecell[{{p{\linewidth}}}]{
                \texttt{\tiny{[GM$\rangle$P2]}}
                \texttt{Mesajın diğer oyuncuya iletildi.} \\
\\ 
\\ 
\texttt{Şu anda en fazla 28 turdan 28. turu oynuyorsunuz.} \\
\texttt{Toplamda 16 cezanın 18 kadarını biriktirdiniz.Diğer oyuncu şu mesajı gönderdi:} \\
\\ 
\\ 
\texttt{"S(2,6) ve M(5,5) olsun. Hemen dene!"} \\
\\ 
\\ 
\texttt{Bir sonraki komutun nedir?} \\
            }
        }
    }
    & \\ \\

    \theutterance \stepcounter{utterance}  
    & & & \multicolumn{4}{p{0.6\linewidth}}{
        \cellcolor[rgb]{0.9,0.9,0.9}{
            \makecell[{{p{\linewidth}}}]{
                \texttt{\tiny{[GM$\langle$P2]}}
                \texttt{HAMLE: S, (2,6)} \\
            }
        }
    }
     \\ \\

    \theutterance \stepcounter{utterance}  
    & & & \multicolumn{2}{p{0.3\linewidth}}{
        \cellcolor[rgb]{0.9,0.9,0.9}{
            \makecell[{{p{\linewidth}}}]{
                \texttt{\tiny{[GM$|$GM]}}
                \texttt{HAMLE: S, (2,6)} \\
            }
        }
    }
    & & \\ \\

    \theutterance \stepcounter{utterance}  
    & & & \multicolumn{2}{p{0.3\linewidth}}{
        \cellcolor[rgb]{0.9,0.9,0.9}{
            \makecell[{{p{\linewidth}}}]{
                \texttt{\tiny{[GM$|$GM]}}
                \texttt{Ceza: (2,6) boş değil; "┼" içeriyor.} \\
\\ 
\texttt{Toplamda 16 cezanın 19 kadarını biriktirdiniz.Lütfen tekrar deneyin!} \\
            }
        }
    }
    & & \\ \\

    \theutterance \stepcounter{utterance}  
    & & \multicolumn{4}{p{0.6\linewidth}}{
        \cellcolor[rgb]{0.9,0.9,0.9}{
            \makecell[{{p{\linewidth}}}]{
                \texttt{\tiny{[GM$\rangle$P2]}}
                \texttt{Ceza: (2,6) boş değil; "┼" içeriyor.} \\
\\ 
\texttt{Toplamda 16 cezanın 19 kadarını biriktirdiniz.Lütfen tekrar deneyin!} \\
            }
        }
    }
    & \\ \\

    \theutterance \stepcounter{utterance}  
    & & & \multicolumn{4}{p{0.6\linewidth}}{
        \cellcolor[rgb]{0.9,0.9,0.9}{
            \makecell[{{p{\linewidth}}}]{
                \texttt{\tiny{[GM$\langle$P2]}}
                \texttt{HAMLE: S, (5,2)} \\
            }
        }
    }
     \\ \\

    \theutterance \stepcounter{utterance}  
    & & & \multicolumn{2}{p{0.3\linewidth}}{
        \cellcolor[rgb]{0.9,0.9,0.9}{
            \makecell[{{p{\linewidth}}}]{
                \texttt{\tiny{[GM$|$GM]}}
                \texttt{HAMLE: S, (5,2)} \\
            }
        }
    }
    & & \\ \\

    \theutterance \stepcounter{utterance}  
    & & & \multicolumn{2}{p{0.3\linewidth}}{
        \cellcolor[rgb]{0.9,0.9,0.9}{
            \makecell[{{p{\linewidth}}}]{
                \texttt{\tiny{[GM$|$GM]}}
                \texttt{Ceza: (5,2) boş değil; "┴" içeriyor.} \\
\\ 
\texttt{Toplamda 16 cezanın 20 kadarını biriktirdiniz.Lütfen tekrar deneyin!} \\
            }
        }
    }
    & & \\ \\

    \theutterance \stepcounter{utterance}  
    & & \multicolumn{4}{p{0.6\linewidth}}{
        \cellcolor[rgb]{0.9,0.9,0.9}{
            \makecell[{{p{\linewidth}}}]{
                \texttt{\tiny{[GM$\rangle$P2]}}
                \texttt{Ceza: (5,2) boş değil; "┴" içeriyor.} \\
\\ 
\texttt{Toplamda 16 cezanın 20 kadarını biriktirdiniz.Lütfen tekrar deneyin!} \\
            }
        }
    }
    & \\ \\

    \theutterance \stepcounter{utterance}  
    & & & \multicolumn{4}{p{0.6\linewidth}}{
        \cellcolor[rgb]{0.9,0.9,0.9}{
            \makecell[{{p{\linewidth}}}]{
                \texttt{\tiny{[GM$\langle$P2]}}
                \texttt{SÖYLE: (2,6) ve (5,2) bende boş değil. M(5,5) dene, S için de (6,7) veya (7,7) dene. Son turdayız, oyunu bitirmek için bitti mi diyelim mi? Yoksa son bir hamle deneyelim mi?} \\
            }
        }
    }
     \\ \\

    \theutterance \stepcounter{utterance}  
    & & & \multicolumn{2}{p{0.3\linewidth}}{
        \cellcolor[rgb]{0.9,0.9,0.9}{
            \makecell[{{p{\linewidth}}}]{
                \texttt{\tiny{[GM$|$GM]}}
                \texttt{SÖYLE: (2,6) ve (5,2) bende boş değil. M(5,5) dene, S için de (6,7) veya (7,7) dene. Son turdayız, oyunu bitirmek için bitti mi diyelim mi? Yoksa son bir hamle deneyelim mi?} \\
            }
        }
    }
    & & \\ \\

    \theutterance \stepcounter{utterance}  
    & & & \multicolumn{2}{p{0.3\linewidth}}{
        \cellcolor[rgb]{0.9,0.9,0.9}{
            \makecell[{{p{\linewidth}}}]{
                \texttt{\tiny{[GM$|$GM]}}
                \texttt{Maximum number of rounds reached} \\
            }
        }
    }
    & & \\ \\

    \theutterance \stepcounter{utterance}  
    & & & \multicolumn{2}{p{0.3\linewidth}}{
        \cellcolor[rgb]{0.9,0.9,0.9}{
            \makecell[{{p{\linewidth}}}]{
                \texttt{\tiny{[GM$|$GM]}}
                \texttt{Initial states:} \\
\texttt{```} \\
\texttt{Player 1:} \\
\texttt{ 1234567} \\
\texttt{╔═╤══╤══╗} \\
\texttt{║◌│◌T│◌◌║ 1} \\
\texttt{║◌└┬─┴──╢ 2} \\
\texttt{║◌◌│O◌◌◌║ 3} \\
\texttt{╟─┬┘P◌┌─╢ 4} \\
\texttt{║◌│AD◌│◌║ 5} \\
\texttt{╟─┼─┬─┤M║ 6} \\
\texttt{║◌│S│◌│◌║ 7} \\
\texttt{╚═╧═╧═╧═╝} \\
\\ 
\\ 
\texttt{Player 2:} \\
\texttt{ 1234567} \\
\texttt{╔═╤══╤══╗} \\
\texttt{║◌│O◌│◌◌║ 1} \\
\texttt{║◌└┬─┴──╢ 2} \\
\texttt{║◌P│◌D◌◌║ 3} \\
\texttt{╟─┬┘◌◌┌─╢ 4} \\
\texttt{║◌│◌◌◌│A║ 5} \\
\texttt{╟─┼─┬─┤S║ 6} \\
\texttt{║M│◌│◌│T║ 7} \\
\texttt{╚═╧═╧═╧═╝} \\
\\ 
\texttt{```} \\
            }
        }
    }
    & & \\ \\

    \theutterance \stepcounter{utterance}  
    & & & \multicolumn{2}{p{0.3\linewidth}}{
        \cellcolor[rgb]{0.9,0.9,0.9}{
            \makecell[{{p{\linewidth}}}]{
                \texttt{\tiny{[GM$|$GM]}}
                \texttt{End states:} \\
\texttt{```} \\
\texttt{Player 1:} \\
\texttt{ 1234567} \\
\texttt{╔═╤══╤══╗} \\
\texttt{║◌│◌◌│◌◌║ 1} \\
\texttt{║◌└┬─┴──╢ 2} \\
\texttt{║AD│O◌◌T║ 3} \\
\texttt{╟─┬┘◌P┌─╢ 4} \\
\texttt{║◌│◌◌◌│◌║ 5} \\
\texttt{╟─┼─┬─┤M║ 6} \\
\texttt{║◌│S│◌│◌║ 7} \\
\texttt{╚═╧═╧═╧═╝} \\
\\ 
\\ 
\texttt{Player 2:} \\
\texttt{ 1234567} \\
\texttt{╔═╤══╤══╗} \\
\texttt{║◌│O◌│◌◌║ 1} \\
\texttt{║◌└┬─┴──╢ 2} \\
\texttt{║AD│◌◌◌T║ 3} \\
\texttt{╟─┬┘◌P┌─╢ 4} \\
\texttt{║◌│◌◌◌│◌║ 5} \\
\texttt{╟─┼─┬─┤S║ 6} \\
\texttt{║M│◌│◌│◌║ 7} \\
\texttt{╚═╧═╧═╧═╝} \\
\\ 
\texttt{```} \\
            }
        }
    }
    & & \\ \\

    \theutterance \stepcounter{utterance}  
    & & & \multicolumn{2}{p{0.3\linewidth}}{
        \cellcolor[rgb]{0.9,0.9,0.9}{
            \makecell[{{p{\linewidth}}}]{
                \texttt{\tiny{[GM$|$GM]}}
                \texttt{* success: True} \\
\texttt{* lose: False} \\
\texttt{* aborted: False} \\
\texttt{{-}{-}{-}{-}{-}{-}{-}} \\
\texttt{* Shifts: 3.00} \\
\texttt{* Max Shifts: 12.00} \\
\texttt{* Min Shifts: 6.00} \\
\texttt{* End Distance Sum: 12.44} \\
\texttt{* Init Distance Sum: 27.62} \\
\texttt{* Expected Distance Sum: 29.33} \\
\texttt{* Penalties: 20.00} \\
\texttt{* Max Penalties: 16.00} \\
\texttt{* Rounds: 28.00} \\
\texttt{* Max Rounds: 28.00} \\
\texttt{* Object Count: 7.00} \\
            }
        }
    }
    & & \\ \\

\end{supertabular}
}

\end{document}
